\chapter{场景编辑与操控}

\paragraph*{本章阅读锚点}
3DGS 编辑的核心矛盾是“显式可改”和“多视角一致”之间的拉扯。原始 3DGS 把场景拆成可直接移动、删改、重着色的高斯，这比 NeRF 更适合交互；但每个高斯同时影响多个视角，若只在单张渲染图上做 2D 编辑，反投回 3D 后容易出现漂浮、遮挡错误和视角不一致。本章按编辑对象来读：风格化方法主要改颜色、SH 和局部几何；文本/语义编辑方法先定位要改的高斯，再用扩散或多视角修补生成监督；交互式方法把用户的点击、拖拽、涂色映射到高斯参数；几何操控方法借 mesh 或 cage 代理传播形变；修复方法则关心删除后背景能否一致补全，以及残留语义是否真正消失。

\section{风格化与艺术编辑}

风格化路线的基础锚点是 StylizedNeRF、ARF 和早期 StyleGaussian：给定一个参考风格图像或文本风格，让同一 3D 场景在所有视角中呈现一致的艺术外观。3DGS 带来的改变是可以直接优化高斯颜色、球谐系数、不透明度、尺度甚至局部几何；困难则是，若只优化颜色，笔触和结构风格不足，若同时改几何，又容易破坏原场景内容。因此这些方法的差别主要在于风格监督来自 VGG/CLIP/扩散模型还是前馈网络，风格作用在颜色、几何还是时空特征上。



% METHOD: 2404.05220
\paragraph{StylizedGS（可控 3DGS 风格化）（2024.04）}
StylizedGS~\cite{2404_05220} 是 3DGS 风格化的早期局部锚点。它相对原始 3DGS 的改变是把风格迁移直接写成高斯参数优化：既调颜色，也允许少量几何和不透明度变化来承载笔触，但用深度保持损失防止几何内容被破坏。

方法分两阶段。第一阶段做全局场景重着色，用颜色直方图匹配让 3DGS 的颜色分布接近风格图像，同时通过基于尺度和不透明度的 filter refinement 去除原重建中的 floaters，避免风格把漂浮噪声染得更明显。第二阶段用最近邻特征匹配损失迁移局部纹理和图案，把 VGG 特征空间中最相近的风格 patch 对齐到渲染图；深度保持损失和正则项约束高斯位置、尺度和不透明度。它还支持颜色、笔触尺度和空间区域控制。公开 StylizedGS README 与论文一致，说明代码实现的是这一两阶段优化流程。
\textbf{实现细节：} 在 DragGaussian 的基础上，DYG（Drag Your Gaussian）（2025.01） 首次实现基于控制点拖拽的 3DGS 几何编辑。前述笼式与网格代理方法均依赖额外的代理结构来传递变形，DYG 的核心改进是去掉代理、直接在 3DGS 上拖拽，并通过隐式三平面表示弥补显式高斯的固有缺陷。其关键创新是多分辨率三平面位置编码器（Multi-resolution Triplane Positional Encoder）：将 3D 空间投影到三个正交平面 $(xy, xz, yz)$ 上，每个平面在多个分辨率下维护可学习特征图，任意位置 $p$ 的特征为三平面多尺度拼接 其中 $\pi_c$ 将点投影到第 $c$ 个平面，$\psi^s$ 在分辨率 $s$ 下进行双线性插值，$\odot$ 为 Hadamard 乘积，$\phi$ 为轻量级融合 MLP。这一设计解决了 3DGS 拖拽编辑中的核心困难——当用户将局部区域拖向新位置时，目标区域往往没有足够的高斯基元（稀疏），导致拖拽后的高斯无法稳定锚定；三平面作为连续隐式几何支架，为目标区域的空白位置提供合理的特征编码，使高斯可以在任意位置获得可靠的位移预测。DYG 还设计了区域感知位置解码器：将高斯分为掩码内集合 $G_m$ 和掩码外集合 $G_{um}$，用两个 MLP 分别预测位移——掩码内高斯正常预测 $\Delta p = N_1(\mathbf{f})$，掩码外高斯通过残差修正 $\Delta p = N_1(\mathbf{f}) + N_2(\mathrm{sg}(\mathbf{f}))$（$\mathrm{sg}$ 为停止梯度算子）避免意外偏移，并辅以正则化损失 $L_\mathrm{reg} = \sum_{g_i \in G_{um}} \|\Delta p_i\|$ 约束非编辑区域不动。整个拖拽过程分为两个阶段：第一阶段冻结所有高斯参数，仅优化三平面编码器和解码器，建立几何支架；第二阶段解冻高斯的颜色、不透明度、旋转和缩放属性并重新启用致密化操作，精化纹理细节。
\textbf{方法信息：} 代码：\url{https://github.com/kristen-z/stylizedgs}。
% METHOD: 2406.08488
\paragraph{ICE-G（图像条件的局部编辑与风格迁移）（2024.06）}
ICE-G~\cite{2406_08488} 的输入是单张参考编辑图像。用户可以手工在参考图中改颜色或纹理，也可以提供风格图；系统需要找到 3D 场景中语义对应区域并迁移变化。它的改变是把 SAM 区域分割和 DINO 特征匹配放到 3DGS 编辑前端。

具体流程先对参考编辑图和训练视图运行 SAM，得到候选区域，再用 DINO 特征为每个训练视图区域寻找参考图中的最佳语义对应。风格迁移分两轮：第一轮用最近邻特征匹配损失传纹理，保证同一语义区域在多视角中一致；第二轮用 L1/SSIM 损失迁移颜色，使输出更鲜艳和稳定。由于底层是可微 3DGS rasterizer，损失直接反传到高斯参数。相对 StyleSplat 需要先建立对象级 3D grouping，ICE-G 更偏“从一张编辑参考图自动找对应区域”。
\textbf{方法信息：} Google / DeepMind 参与了这篇论文。
\textbf{设计动机与边界：} 相对原始 3DGS 以独立高斯和图像重建误差为主的设定，ICE-G 把问题落在以下可实现的改动上：ICE-G支持两种编辑界面：手动纹理/颜色编辑和从示例图像自动迁移。核心流程：(1) 使用SAM对编辑图像和采样数据集视图进行分割，生成掩码区域。(2) 利用DINO特征在每个采样视图的每个区域找到编辑图像中最佳语义对应区域，通过最近邻特征匹配（NNFM）建立对应关系。(3) 风格迁移分两轮：第一轮用NNFM损失迁移纹理（保证3D一致性），第二轮用L1/SSIM损失迁移颜色（更鲜艳的颜色效果）。由于高斯泼溅的可微分光栅化器，编辑可直接反向传播到高斯参数。方法也适用于NeRF但速度较慢。其设计目的不是单独增加一个模块，而是用上述表示、监督或渲染约束消除原流程在目标任务中的具体失配。提出基于单张参考图像快速编辑3D模型的方法。利用SAM分割编辑图像，通过DINO特征匹配语义对应区域，自动将颜色/纹理变化应用到对应视图。支持手动局部编辑、基于对应的风格迁移和多种风格组合。基于高斯泼溅实现快速训练和局部编辑。适用时还应保留论文给出的边界：仅支持颜色和纹理编辑，不支持形状修改；风格迁移依赖DINO特征匹配的质量。
% METHOD: 2407.09473
\paragraph{StyleSplat（对象级 3D 风格迁移）（2024.07）}
StyleSplat~\cite{2407_09473} 把 StylizedGS 的全局风格化推进到对象级。它的问题起点是：场景里可能只有一个物体需要变成某种材质或画风，其他物体和背景不应改变。直接用 2D 掩码约束单视图风格会在其他视角漏改或误改，因此方法先做 3D 高斯分组。

具体流程先用 DEVA 和 SAM 生成多视角、时序一致的 2D 分割掩码。然后借鉴 Gaussian Grouping，为每个高斯附加紧凑特征向量，通过可微渲染把特征投到图像上，并用 2D 掩码监督高斯特征聚类。用户选中某个对象后，方法只优化该对象高斯的球谐系数或颜色参数，用最近邻特征匹配或 Gram 风格损失与参考风格对齐。相对 ReCoGS 的区域重着色，StyleSplat 更强调对象选择和风格纹理迁移。
\textbf{实现细节：} 与 StylizedGS 的场景级全局风格化不同，StyleSplat（2024.07） 将风格化粒度推进到物体级——场景中的不同物体可以被赋予不同的艺术风格，而这正是场景级方法无法实现的核心局限。动机。场景级方法对所有区域施加统一风格，无法让场景中不同物体呈现不同艺术效果；StyleSplat 通过 Gaussian Grouping 机制先分割再独立风格化，实现场景内多物体独立风格化。其关键设计是 Gaussian Grouping 机制：为每个 3D 高斯分配一个 16 维可学习特征向量 $\mathbf{e}_i \in \mathbb{R}^{16}$，该编码与球谐系数同步优化，渲染时通过 $\alpha$-混合聚合为像素级特征： 渲染特征经分类器输出语义标签，以交叉熵损失与 DEVA+SAM 生成的时序一致 2D 分割掩码对齐，辅以空间一致性损失确保近邻高斯特征相似。分割完成后，仅解冻用户指定物体的球谐系数进行梯度更新，其余高斯属性全部冻结，风格化被严格约束在目标物体内部。NNFM 损失仅对选物体的 VGG 特征计算，实现首个支持场景内多物体独立风格化的轻量方法。
\textbf{方法信息：} 代码：\url{https://github.com/bernard0047/style-splat}。
\textbf{设计动机与边界：} 相对原始 3DGS 以独立高斯和图像重建误差为主的设定，StyleSplat 把问题落在以下可实现的改动上：StyleSplat的核心流程分三步。第一步：使用DEVA（解耦视频分割方法）结合SAM（分割一切模型）生成时序一致的2D分割掩码，每个掩码对应场景中的一个物体。第二步：基于Gaussian Grouping方法，每个3D高斯附加16维可学习紧凑特征向量，类似球谐系数的优化方式。渲染时通过alpha blending聚合特征得到像素级特征图，使用2D掩码作为监督训练特征聚类。第三步：风格迁移阶段，用户选择目标物体并提供参考风格图像，使用最近邻特征匹配损失将选定高斯的球谐系数（最高4阶）与风格图像对齐。具体通过VGG特征空间的Gram矩阵匹配实现风格迁移，同时保持场景几何和其他物体不变。整个流程无需训练额外网络，直接微调现有高斯参数。其设计目的不是单独增加一个模块，而是用上述表示、监督或渲染约束消除原流程在目标任务中的具体失配。提出首个支持3DGS场景中局部物体级风格迁移的轻量方法。通过2D掩码生成和3D高斯分割实现精确物体选择，利用最近邻特征匹配损失确保风格一致性，支持同一场景中多个物体应用不同风格。适用时还应保留论文给出的边界：风格质量依赖参考图像的质量和风格复杂度；分割精度受限于SAM的2D分割能力；大幅视角变化下风格一致性可能下降
% METHOD: 2410.10412
\paragraph{4DStyleGaussian（零样本 4D 风格迁移）（2024.10）}
4DStyleGaussian~\cite{2410_10412} 把风格化从静态 3D 推到动态 4D。它继承 4DGS 的时空高斯表示，但普通 3D 风格损失无法保证同一动态对象在不同时间步保持相同风格；逐帧风格化又会产生明显闪烁。

方法在 4DGS 训练中为每个高斯学习嵌入特征，通过可微 splatting 渲染出嵌入特征图，再用 MLP 得到颜色图和特征图。内容特征由 2D 可逆神经网络提取，减少风格变换时的信息丢失；风格侧预测 4D 风格变换矩阵，使同一个高斯在不同时间和视角下共享时空一致的风格映射。相对 Style4D-Bench 的基准式训练，4DStyleGaussian 更强调零样本和实时任意风格迁移。
\textbf{方法信息：} 代码：\url{https://github.com/liangwanlin/4dstylegaussian}。
\textbf{设计动机与边界：} 相对原始 3DGS 以独立高斯和图像重建误差为主的设定，4DStyleGaussian 把问题落在以下可实现的改动上：方法基于4D-GS骨干，包含：（1）嵌入4DGS训练：为每个高斯学习嵌入特征，通过可微splatting和光栅化获得嵌入特征图，分别通过两个MLP获取颜色图和特征图。使用2D可逆神经网络从内容图像提取特征，避免内容保真度损失。（2）4D风格变换矩阵预测：利用嵌入的4D高斯预测时空一致的风格变换矩阵，实现跨视角和时间的风格一致性迁移。可逆神经网络确保特征提取过程中不丢失内容信息。其设计目的不是单独增加一个模块，而是用上述表示、监督或渲染约束消除原流程在目标任务中的具体失配。提出4DStyleGaussian，首个零样本4D风格迁移框架，通过可逆神经网络和4D风格变换矩阵实现实时任意风格迁移，同时保持内容保真度、多视角一致性和时序连贯性。适用时还应保留论文给出的边界：风格迁移质量受限于预训练编码器；对极端风格变换可能产生伪影；训练需要4D视频数据
% METHOD: 2412.03371
\paragraph{SGSST（超高分辨率多尺度风格迁移）（2024.12）}
SGSST~\cite{2412_03371} 关注的是超高分辨率 3DGS 场景。普通风格损失在 4K 级渲染中要么只匹配整体调色板，要么分尺度独立优化，难以同时保留大色块和近距离笔触。SGSST 的改变是提出 SOS，即 Simultaneously Optimized Scales，用一个无参数、可解释的多尺度损失同时约束所有尺度。

方法直接优化预训练 3DGS 的属性参数，例如颜色和球谐系数。SOS 基于全局神经统计，把不同尺度的风格特征放到同一优化式中，使远看时整场景保留画作整体调色板，近看时仍能看到笔触或纹理细节。实现上先按原始 3DGS 流程训练真实场景，再运行 SGSST 的 stylization 脚本；仓库还要求克隆 Scaling Painting Style Transfer 作为 SOS 风格损失依赖。相对 StylizedGS 的两阶段颜色与细节优化，SGSST 的重点是多尺度统一损失，而不是新增交互式编辑接口。
\textbf{实现细节：} 前馈与零样本风格化沿"效率提升"与"模态扩展"推进。InstantStyleGaussian 简化流程实现高效迁移，4DStyleGaussian 与 ZDySS 将零样本扩展到动态 4D 场景；架构上 Stylos 用双路径 Transformer 实现单次前向多视角风格化，AnyStyle 与 M2StyleGS 引入多模态参考，Multi-StyleGS 支持同场景多风格切换。对齐、一致性与解耦控制是另一核心张力——ABC-GS 通过分割掩码实现可控对齐，Reference-based 方法 以参考图引导可控风格化；SSGaussian 同时保证语义感知与结构保持，SGSST 扩展到超高分辨率；艺术化路径上 ArtNVG 分离内容与风格，StyleMe3D 用多编码器解耦先验提升三维保真度，Morpheus 联合编辑形状与颜色。效率与基准方面，免优化方法 绕过迭代实现即时迁移，FantasyStyle 通过可控蒸馏提升质量，Style4D-Bench 提供 4D 评测基准，零样本增强方法 利用 2D 预训练先验；领域扩展中 TexGS-VolVis 进入体积可视化，3D 多视角风格化方法 通过无姿态对应匹配鲁棒保持几何；特定风格类型上 AI 驱动方法 扩展到 3D 环境，GT2-GS 利用几何感知纹理转移，Thinking Like Van Gogh 以流引导实现结构感知，DiffStyle3D 通过注意力优化保证一致性，Tune-Your-Style 提供强度可调迁移。
\textbf{方法信息：} 代码：\url{https://github.com/jianlingwang2021/sgsst}。
\textbf{设计动机与边界：} 相对原始 3DGS 以独立高斯和图像重建误差为主的设定，SGSST 把问题落在以下可实现的改动上：SGSST是一种基于优化的方法，将风格迁移应用到预训练的3DGS场景。核心创新是SOS(Simultaneously Optimized Scales)损失，该损失基于全局神经统计，以单一无参数可解释公式实现多尺度风格迁移。与传统需要分尺度独立优化的方法不同，SOS在单次优化中同时处理所有尺度。这使得当用户在3D环境中靠近物体时，能看到精细的绘画纹理(如笔触)，而全局场景保持风格画作的整体色彩调色板。方法直接优化3DGS的属性参数(球谐系数等)。其设计目的不是单独增加一个模块，而是用上述表示、监督或渲染约束消除原流程在目标任务中的具体失配。提出SOS(Simultaneously Optimized Scales)损失，一种无参数、可解释的单公式多尺度损失，首次实现超高分辨率(4K以上)3DGS场景风格迁移。通过全局神经统计实现多尺度纹理(如画笔笔触)的精确传递，在视觉质量上显著优于现有方法。适用时还应保留论文给出的边界：基于优化的方法需要逐场景优化耗时较长；对极端抽象艺术风格的效果可能有限；未探索实时风格切换。
% METHOD: 2505.22854
\paragraph{CLIPGaussian（统一多模态风格插件）（2025.05）}
CLIPGaussian~\cite{2505_22854} 把风格迁移扩展到 2D、视频、3D 和 4D 场景。它的锚点是 CLIP 引导的语义风格对齐，而不是单一 VGG Gram loss。普通 3DGS 风格化往往只支持参考图像，或者只改静态场景；CLIPGaussian 的改变是用同一插件式模块处理文本和图像条件，并在 3D/4D 设置中联合优化颜色和几何。

具体做法是在已有 Gaussian Splatting 流程中加入 CLIP 风格目标：文本提示或参考图像被编码到 CLIP 特征空间，渲染出的图像也被编码并与目标风格对齐。对于 3D 和 4D 高斯，方法不仅改颜色，还允许对几何参数做受约束调整，使风格从“换色”扩展到形状和纹理气质。动态场景中，它通过共享潜在表示和一致性约束保持跨帧风格稳定。相对 4DStyleGaussian 的专门动态风格矩阵，CLIPGaussian 更强调跨模态、跨任务的统一插件。
\textbf{实现细节：} 动机。风格迁移方法高度碎片化，不同模态和任务需要不同工具；CLIPGaussian 的核心洞察是：所有模态都可以用高斯基元统一表示，一旦统一，风格迁移就归结为同一操作——微调高斯参数。直觉上，CLIPGaussian 就像是给不同语言的翻译器装上同一个引擎——一旦所有内容都统一为高斯基元表示，风格迁移就变成同一套操作流程。CLIPGaussian（2025.05） 的关键洞察是：尽管这些模态的数据形式不同，它们都可以用高斯基元统一表示——2D 图像用 MiRaGe 的 2D 高斯、视频用 VeGaS 的时序高斯、3D 场景用经典 3DGS、4D 场景用 D-MiSo 的动态高斯。一旦统一为高斯表示，风格迁移就可以归结为同一操作：微调高斯参数。具体而言，CLIPGaussian 采用两阶段流程。第一阶段用现有高斯模型训练场景，得到基座模型 $\mathcal{G}$。第二阶段冻结高斯数量，仅微调颜色（SH 系数）和几何参数，通过四项损失的加权组合驱动风格化： 其中方向 CLIP 损失 $\mathcal{L}_d$ 在全局尺度上将风格因子 $S$（文本提示或风格图像）的 CLIP 方向与重建图像的 CLIP 方向对齐；补丁 CLIP 损失 $\mathcal{L}_p$ 对随机裁剪的局部块施加同样的方向约束，确保风格细节渗透到局部纹理而非仅改变全局色调；VGG 内容损失 $\mathcal{L}_c$ 保持风格化结果与原始内容在 conv4\_2/conv5\_2 特征上的相似度；背景损失 $\mathcal{L}_b$ 通过 L1 距离约束背景像素不引入风格伪影。由于不改变高斯数量，风格化后的模型与原始模型体积相同，且支持通过线性插值高斯参数实现风格之间的平滑过渡。整个方法作为插件模块集成到任意高斯流程中，无需大型生成模型或从头重新训练。
\textbf{方法信息：} 代码：\url{https://github.com/kornelhowil/clipgaussian}。
\textbf{设计动机与边界：} 相对原始 3DGS 以独立高斯和图像重建误差为主的设定，CLIPGaussian 把问题落在以下可实现的改动上：CLIPGaussian的核心设计理念是直接在基于高斯溅射模型的架构内操作，作为插件类型组件集成到现有流程：1)统一框架：支持文本提示和参考图像两种条件输入，覆盖2D图像、视频、3D物体和4D动态场景四种数据模态；2)插件式设计：集成到现有GS流程无需重新训练基础模型，利用基础模型的架构进行端到端优化所有参数，不改变模型大小；3)颜色-几何联合优化：在3D和4D设置中，同时优化颜色和几何参数，实现超越简单颜色变化的风格迁移；4)时间一致性：在视频和4D场景中，通过共享潜在表示和一致性约束确保跨帧风格的一致性应用；5)CLIP引导：利用CLIP模型的语义特征空间引导风格迁移过程，确保风格在语义层面的一致应用。其设计目的不是单独增加一个模块，而是用上述表示、监督或渲染约束消除原流程在目标任务中的具体失配。提出CLIPGaussian，首个统一风格迁移框架，支持文本和图像引导的风格化，覆盖2D图像、视频、3D物体和4D场景四种模态。方法直接在高斯基元上操作，作为插件模块集成到现有GS流程中，无需大型生成模型或从头重新训练。在3D和4D设置中支持颜色和几何的联合优化，在视频中实现时间一致性，同时保持模型大小不变。展示了跨所有任务的优越风格保真度和一致性。适用时还应保留论文给出的边界：风格迁移质量受限于基础GS模型的表示能力；复杂风格可能难以精确再现；4D场景的计算开销较大。
% METHOD: 2508.05813
\paragraph{Optimization-Free Style Transfer（直接处理 splat 文件的无优化风格迁移）（2025.08）}
Optimization-Free Style Transfer for 3D Gaussian Splats~\cite{2508_05813} 选择了一条很实用的路线：用户只有一个公开的 \texttt{.ply} 或 \texttt{.splat} 文件，没有原始训练相机和图像，也不想重新优化。相对 StylizedGS 这类需要渲染训练视图并反传到高斯参数的基线，该方法把高斯中心转成图结构，把问题改成图上风格传播。

具体做法以每个高斯中心为图节点，使用 Open3D 和 Ball-Pivoting 估计局部法线，通过 $K=16$ 近邻连边，并用局部平面投影确定边方向。构图后，方法把 Interpolated SelectionConv 用作表面卷积，在图上运行预训练图像风格迁移网络，把风格化颜色分配给节点，再插值回原始高斯属性。它不改几何、不需要相机、不做梯度下降，因此速度快、门槛低；代价是风格主要是观察者无关的颜色/纹理迁移，难以像 StylizedGS 那样联合几何笔触。
\textbf{实现细节：} 语言驱动解耦：3DitScene。3DitScene（2024.05） 保留了 Point'n Move 的对象级操控能力，但将分割方式从"用户点击 → SAM 分割"替换为"语言/框选 → CLIP 语义匹配"，使用户无需精确点击即可选中目标。具体地，3DitScene 为每个 3D 高斯基元额外嵌入一个 CLIP 语言向量 $e \in \mathbb{R}^C$，利用 3DGS 的可微分泼溅将该向量渲染为 2D 语义特征图 $E \in \mathbb{R}^{C \times H \times W}$，以蒸馏损失。
\textbf{方法信息：} 代码：\url{https://github.com/davidmhart/fastsplatstyler}。
\textbf{设计动机与边界：} 相对原始 3DGS 以独立高斯和图像重建误差为主的设定，Optimization-Free Style Transfer for 3D Gaussian Splats 把问题落在以下可实现的改动上：该方法基于Interpolated SelectionConv框架，将3DGS表示转化为图结构以支持标准图像卷积操作。具体流程包括：(1)点云采样：将每个高斯的中心作为图节点；(2)法线估计：使用Ball-Pivoting方法通过Open3D估计法线向量；(3)KNN边构建：使用K=16的K近邻方法连接节点；(4)平面投影选择：通过局部平面投影确定边的方向。构建好图结构后，使用预训练的图像风格迁移网络（Li等人的方法）对图上节点进行风格化处理，然后将风格化颜色通过插值映射回原始高斯。该方法不需要原始训练图像，可以直接处理公开的.ply或.splat文件，风格化结果保存在高斯属性中与观察者无关。其设计目的不是单独增加一个模块，而是用上述表示、监督或渲染约束消除原流程在目标任务中的具体失配。提出首个无需重建、无需优化、无需原始相机视角的3DGS风格迁移方法。通过在3DGS隐式表面上构建图结构，利用SelectionConv表面卷积网络和预训练图像风格迁移网络直接对.ply或.splat文件进行风格化处理，速度极快（CPU上不到2分钟）。适用时还应保留论文给出的边界：风格化质量受限于图结构对表面的近似精度；对于体积分布较散的高斯表示效果可能下降；风格化结果的分辨率受图密度限制。
% METHOD: 2508.08136
\paragraph{FantasyStyle（扩散蒸馏与频率一致性）（2025.08）}
FantasyStyle~\cite{2508_08136} 的锚点是扩散模型驱动的 3D 风格迁移。它指出 VGG 风格损失容易把参考图像中的具体内容泄漏到场景里，而直接用 SDS/DDS 又会带来跨视角不一致和过强重建约束。FantasyStyle 的改变是把多视角一致性放到扩散隐空间的频率分解中处理。

方法先对多视角渲染图做 DDIM 噪声化，在隐变量上执行 3D FFT，分析不同频段的跨视角稳定性。它发现低频分量更容易携带视角相关局部差异，高频纹理反而更稳定地代表风格，因此通过 Multi-View Frequency Consistency 选择性衰减低频、保留纹理高频。随后在可控风格蒸馏中引入 negative guidance，排除风格图像里不希望迁移的内容，并弱化重建项，主要微调 3DGS 颜色参数。相对 StylizedGS 的最近邻风格损失，它更依赖扩散先验；相对 AnyStyle，它仍是按场景优化而非一次前馈。
\textbf{实现细节：} DGE（Direct Gaussian Editor）（2024.04） 正是针对这一瓶颈提出了替代思路：与其在 3D 空间反复修补不一致的编辑结果，不如直接在图像空间保证多视角一致性，然后仅需将 3DGS 拟合到编辑后的图像即可完成 3D 编辑。与 GaussianEditor 通过 HGS 逐代凝固高斯属性来抑制迭代发散不同，DGE 将核心难题前移至 2D 编辑阶段——借鉴视频编辑中无需额外训练即可扩展图像生成器的思路，将多视角序列视为轨道视频，使多视角一致性等价于视频编辑中的时序一致性。DGE 将流程分为两步：关键视角编辑和特征注入。关键视角编辑阶段，在 IP2P 的自注意力层中引入时空注意力，使 $T$ 个视角的 query 同时关注所有视角的 key-value： 但该操作计算开销随视角数二次增长，因此实际执行中仅对随机采样的关键视角子集 $\mathcal{K} \subset \{1, \ldots, T\}$ 执行。在特征注入阶段，DGE 利用极线约束（Epipolar Constraint）将关键视角的编辑特征传播到其余视角：给定非关键视角 $t'$ 的特征图 $\mathbf{F}_{t'}$ 和最近关键视角 $k^{*}$ 的特征图 $\mathbf{F}_{k^{*}}$，对应关系通过基础矩阵 $F$ 约束搜索空间获得： 其中 $D$ 为余弦距离，$Fv \parallel u$ 表示点 $v$ 在关键视角中的对应点必须位于极线 $Fu$ 上。这一约束将对应搜索从二维平面压缩到一维极线上，降低了误匹配率。最终 DGE 通过渲染损失 $\theta' = \arg\min_\theta \sum_{t=1}^{T} \| I'_t - R(\theta, \pi_t) \|$ 直接优化 3DGS 参数，整个编辑流程无需额外训练扩散模型。
\textbf{实现补充：} 实现上还要核对输入、输出、参数、特征、初始化、训练损失、正则、渲染或推理接口，以及优化中的剪枝、分裂、采样、编码或解码步骤；这些抓手决定该方法到底改了 3DGS 管线的哪一环。
\textbf{方法信息：} 代码：\url{https://github.com/yangyt46/fantasystyle}。
% METHOD: 2508.19243
\paragraph{Style4D-Bench（动态场景风格化基准与基线）（2025.08）}
Style4D-Bench~\cite{2508_19243} 的贡献一半是 benchmark，一半是基线方法。静态 3D 风格化只需保证多视角一致；4D 风格化还要保证时间一致，否则同一动态物体的笔触会在帧间跳动。它用 4DGS 作为几何锚点，并定义评估空间保真度、时间一致性和多视角一致性的协议。

基线 Style4D 先用 4DGS 表示动态场景几何，再为每个高斯引入轻量逐高斯 MLP，输入时间和深度信息，输出风格化颜色，形成 Style Gaussian Representation。训练目标包含风格保真、对比一致性学习和结构内容保持：相邻时空点在风格特征上应接近，几何结构不能被风格扰动破坏。相对 4DStyleGaussian 的零样本实时迁移，Style4D-Bench 更像给动态风格化建立测量尺和标准基线。
\textbf{实现细节：} 风格化领域还有多种值得关注的新方向，以下从数据基准、特定风格和用户控制三个维度加以梳理。在数据与基准方面，DS-StyleGaussian（2026.05） 利用 2D 预训练解码器弥补 3D 风格迁移的数据稀缺问题，Style4D-Bench（2025.08） 则提出首个 4D 风格化基准，为时序一致性评估提供标准。
\textbf{实现补充：} 为避免只把它读成名称或数据点，还应把实现接口明确拆开：输入是什么，输出是高斯参数、图像、轨迹、标注还是评测指标，哪些参数、网络模块、训练损失或渲染接口真正接到原始 3DGS / Scaffold-GS / 4D-GS 管线上，以及剪枝、分裂、采样、解码、优化或评测协议在哪一步发生。若该条目是survey、benchmark 或 dataset，这里的“实现”对应论文集合、任务划分、相机/轨迹/标注组织、评价指标和基线配置，而不是虚构新的 rasterizer；若是算法条目，则对应高斯位置、尺度、旋转、不透明度、特征或 MLP 解码如何进入训练和渲染流程。
\textbf{方法信息：} Alibaba 参与了这篇论文；代码：\url{https://github.com/Becky-catherine/Style4D-Bench}。



% METHOD: 2509.26455
\paragraph{Stylos（单次前向多视图风格化）（2025.09）}
Stylos~\cite{2509_26455} 与 AnyStyle 同属前馈路线，但它的输入范围更宽，可以从单视图到多视图集合直接预测风格化 3D 高斯和相机。它继承 VGGT 的交替注意力主干，用内容分支推断几何，用风格分支注入外观，避免风格干扰深度、尺度和位姿。

具体结构是双路径 Transformer。内容图像通过帧级和全局自注意力推断位置、深度、尺度、旋转和不透明度；风格参考图像通过 CrossBlock 的全局交叉注意力进入颜色系数预测。为了避免只在 2D 渲染图上匹配风格导致视角不一致，它提出体素级 3D 风格损失：把高斯按空间位置聚合到体素网格，计算体素特征统计，再与风格图像统计对齐。相对 AnyStyle 的轻量控制支路，Stylos 更像完整的风格化重建模型。
\textbf{方法信息：} 代码：\url{https://github.com/hanzhouliu/stylos}。
\textbf{设计动机与边界：} 相对原始 3DGS 以独立高斯和图像重建误差为主的设定，Stylos 把问题落在以下可实现的改动上：Stylos核心设计：(1) 双路径Transformer主干：继承VGGT的交替注意力架构（帧级+全局自注意力）处理内容图像，推理几何属性（位置、深度、尺度、旋转、不透明度）。风格通过全局交叉注意力注入（CrossBlock模块：自注意力+交叉注意力），引导颜色系数预测。几何仅从内容派生，风格仅影响外观。(2) 风格聚合器：将风格参考图像注入Transformer处理流程，通过CrossBlock在每一层将风格特征与内容特征融合。(3) 体素级3D风格损失：将3D高斯按空间位置聚合到体素网格，计算体素级特征统计（均值/方差），与风格图像的Gram矩阵统计对齐。这比2D风格损失更好地保证多视角一致性。(4) 联合预测：同时输出3D高斯参数和相机参数（位姿来自VGGT主干的推理）。支持任意数量输入视图。其设计目的不是单独增加一个模块，而是用上述表示、监督或渲染约束消除原流程在目标任务中的具体失配。提出Stylos，单次前向通过的3D高斯风格化框架，无需逐场景优化或预计算位姿，支持从单视图到多视图的零样本3D风格迁移。设计双路径Transformer主干：几何预测保留自注意力保持几何保真度，风格通过全局交叉注意力注入确保视觉一致性。提出体素级3D风格损失对齐聚合场景特征与风格统计。适用时还应保留论文给出的边界：风格化质量受限于T2I模型的先验知识。在极端几何复杂场景中风格一致性可能降低。对风格参考图像的选择敏感。
% METHOD: 2511.18441
\paragraph{ReCoGS（交互式区域重着色）（2025.11）}
ReCoGS~\cite{2511_18441} 更像艺术工具中的局部重着色，而不是全局风格化。用户在一个视图中画出 2D 选择掩码，系统要把这个选择稳定映射到 3DGS 的高斯集合，再在所有训练视图中一致改变颜色。

具体流程先从当前视图渲染深度，使用基于高斯深度中位数和透射率阈值的启发式方法，把 2D 掩码反投影成 3D 点云。再把这些 3D 点投影回所有训练视图，得到多视角编辑掩码。在掩码区域内执行重新着色操作后，后台持续用光度损失优化场景参数，用户可以继续交互调整。相对 StylizedGS 的参考图像风格损失，ReCoGS 的核心是 2D 选择到 3D 区域的可靠传递。
\textbf{方法信息：} 代码：\url{https://github.com/loryruta/recogs}。
\textbf{设计动机与边界：} 相对原始 3DGS 以独立高斯和图像重建误差为主的设定，ReCoGS 把问题落在以下可实现的改动上：ReCoGS管线包含：用户在当前视点绘制2D选择掩码；通过深度图预测（基于高斯深度中位数截断策略）将2D掩码反投影为3D点云；将3D点云投射到所有训练视图生成2D编辑掩码；在编辑后的训练视图上应用重新着色操作；后台使用光度损失优化场景参数。深度图预测采用启发式方法：当透射率T低于阈值时取最后高斯的z值作为像素深度。场景参数优化持续进行，用户可连续编辑选择掩码。整个流程在交互式编辑工具中实现。其设计目的不是单独增加一个模块，而是用上述表示、监督或渲染约束消除原流程在目标任务中的具体失配。提出ReCoGS交互式重着色管线，允许用户在预训练3DGS场景上精确选择和重新着色区域。通过2D选择掩码反投影为3D点云再投射到训练视图实现精确区域选择，后台优化场景参数实现实时编辑效果。适用时还应保留论文给出的边界：深度图预测依赖启发式方法可能不精确；复杂几何区域的选择可能不准确；仅支持着色编辑
% METHOD: 2602.04043
\paragraph{AnyStyle（多模态前馈风格化）（2026.02）}
AnyStyle~\cite{2602_04043} 是风格化路线中的前馈代表性锚点。它继承 AnySplat 一类 feed-forward 3DGS 重建器，但把输出从“还原输入图像”扩展为“按文本或参考图像风格重建”。传统 StylizedGS 需要对每个场景再优化，FantasyStyle 依赖扩散蒸馏；AnyStyle 的改变是把风格注入做成轻量分支，单次前向即可输出风格化高斯。

具体结构是冻结原始重建骨干以保留几何，另建风格分支复制 Gaussian Head 和 Aggregator 用于调整外观。文本和图像风格都映射到共享 CLIP 嵌入空间，因此可以只用文本、只用图像，也可以在二者之间插值。风格注入使用类似 ControlNet 的 zero-initialized convolution，训练初期不扰动原重建，随后逐步学习如何把风格条件注入颜色预测。它相对 Stylos 的差异在于，Stylos 直接用双路径 Transformer 从内容和风格共同预测高斯；AnyStyle 更像给已有前馈 3DGS 骨干接一条可替换的风格控制支路。
\textbf{实现细节：} 沿前馈路线进一步推进，AnyStyle（2026.02） 基于 ControlNet 启发的零初始化卷积实现架构无关的风格注入，仅需轻量微调即可集成到预训练 3DGS 骨干网络中，支持文本和图像双模态控制。M2StyleGS（2026.04） 首次实现多模态（文本/图像）参考的实时 3DGS 风格迁移，提出 subdivisive flow 精确对齐 CLIP 与 VGG 特征空间。
\textbf{实现补充：} 实现上还要核对输入、输出、参数、特征、初始化、训练损失、正则、渲染或推理接口，以及优化中的剪枝、分裂、采样、编码或解码步骤；这些抓手决定该方法到底改了 3DGS 管线的哪一环。
\textbf{设计动机与边界：} 相对原始 3DGS 以独立高斯和图像重建误差为主的设定，AnyStyle 把问题落在以下可实现的改动上：基于AnySplat骨干网络构建双分支架构：冻结骨干分支（保留几何结构）+风格分支（包含复制的Gaussian Head和Aggregator用于外观调整）。多模态条件机制：文本或图像风格输入通过共享CLIP嵌入空间映射，支持风格图像、文本提示及两者间的连续插值。风格注入核心：受ControlNet启发的零初始化卷积（zero-initialized convolution），在训练初期不影响原始重建，逐渐学习风格化能力。训练策略：从预训练AnySplat初始化，仅添加轻量风格注入模块进行微调，而非从头训练。支持冻结骨干替换（如切换为DepthAnything3），架构无关设计。其设计目的不是单独增加一个模块，而是用上述表示、监督或渲染约束消除原流程在目标任务中的具体失配。提出AnyStyle，首个支持多模态（文本+图像）零样本风格化的前馈3DGS重建框架。基于ControlNet启发的零初始化卷积实现架构无关的风格注入，仅需轻量微调即可集成到预训练3DGS骨干网络中，用户研究确认风格化质量优于SOTA。适用时还应保留论文给出的边界：依赖预训练骨干质量；复杂风格（如大幅几何变化）效果有限；风格注入模块仍需额外训练。
\textbf{方法信息：} Microsoft 参与了这篇论文；代码：\url{https://github.com/joaxkal/AnyStyle}。



% METHOD: 2606.24144
\paragraph{Geometry-Aware Style Transfer（几何感知风格迁移）（2026.06）}
Geometry-Aware Style Transfer in 3D Gaussian Splatting~\cite{2606_24144} 试图补上多数 3DGS 风格化方法的短板：只改颜色和球谐系数时，参考风格中的笔触、轮廓和局部结构很难进入三维场景；但如果把高斯位置、尺度和颜色一起优化，外观变化又容易把几何拉坏。它相对 StylizedGS、SGSST 这一类外观风格化路线的主要改变，是把“颜色像风格”和“几何带有风格结构”拆成两个交替优化的问题。

具体做法是先从已有 3DGS 和一张风格图像出发，用颜色匹配把整体色彩分布拉到风格域，避免后续优化从过大的外观差异开始。随后进入解耦优化：颜色阶段主要更新外观参数，使多视角渲染图在风格特征上接近参考图；几何阶段才有限度地更新位置、尺度或形状，让局部边缘和结构也带上风格图案。两阶段都使用 geometry-aware contrastive feature matching：渲染视图不仅比较 RGB 特征，还引入深度和边缘估计形成几何线索，正样本把内容区域拉向相似风格区域，负样本抑制错误匹配。这样读者可以把它理解为“风格化损失不再只看颜色纹理，而是把可控的局部几何变化也纳入风格迁移”。它不是通用几何编辑器，而是风格化路线中较明确地允许几何参与的一支。
\textbf{实现细节：} Geometry-Aware Style Transfer in 3D Gaussian Splatting。Geometry-Aware Style Transfer in 3D Gaussian Splatting（2026.06） 延续 3DGS stylization 在"风格迁移与几何编辑"上的基本设定，主要改变是：Geometry-Aware Style Transfer：颜色与几何参数解耦交替优化，并用几何适配机制保持稳定。3DGS 的显式高斯表示和可微渲染能力保持为基础。其设计动机是：将3DGS stylization从颜色迁移扩展到结构级几何适配，并通过解耦优化减少颜色和几何更新互相干扰。具体而言，方法交替更新颜色和几何参数，在保持场景一致性的同时把风格中的结构属性迁移到高斯场。适用边界是：几何风格迁移可能破坏真实度；风格和结构解耦程度需要调节。
\textbf{方法信息：} 代码：\url{https://github.com/oweixx/gast}。
\textbf{设计动机与边界：} 相对原始 3DGS 以独立高斯和图像重建误差为主的设定，Geometry-Aware Style Transfer in 3D Gaussian Splatting 把问题落在以下可实现的改动上：方法交替更新颜色和几何参数，在保持场景一致性的同时把风格中的结构属性迁移到高斯场。其设计目的不是单独增加一个模块，而是用上述表示、监督或渲染约束消除原流程在目标任务中的具体失配。将3DGS stylization从颜色迁移扩展到结构级几何适配，并通过解耦优化减少颜色和几何更新互相干扰。适用时还应保留论文给出的边界：几何风格迁移可能破坏真实度；风格和结构解耦程度需要调节。
% METHOD: 2606.26754
\paragraph{Capacity-Controlled Multi-View Stylization（容量约束的多视角风格分配）（2026.06）}
Capacity-Controlled Multi-View Stylization of 3D Gaussian Splatting~\cite{2606_26754} 关注的不是“能不能把风格传进 3DGS”，而是风格图案在多视角中如何分配。普通 2D feature matching 往往逐视图独立运行，容易出现两类问题：一是很多内容 patch 都匹配到同一个风格 patch，导致重复纹理；二是同一片三维表面在不同视角中匹配到不同风格区域，造成闪烁或纹理跳变。它的改变是把局部风格匹配改写为带容量约束的半平衡最优传输问题。

流程上，方法先训练一个更适合细纹理承载的 3DGS，通过几何正则让高斯更小、更均匀，减少大高斯吞掉细节的问题。风格化阶段使用 capacity-controlled feature transport loss：内容视图的局部特征和风格图像的局部特征之间不再做无约束最近邻匹配，而是在最优传输中限制每个风格特征可被复用的容量。直观地说，一块风格纹理不能被无限复制到全场景，风格图案会被更均衡地分配到不同内容区域。为了解决跨视角不一致，方法还加入 cross-view matching guidance，让相邻视角中对应的三维区域倾向于使用一致的风格匹配。相对 GAST 更强调颜色与几何的交替更新，这篇工作的核心在匹配分配和视角一致性，适合作为多视角风格化损失设计的代表。
\textbf{实现细节：} Capacity-Controlled Multi-View Stylization of 3D Gaussian Splatting。Capacity-Controlled Multi-View Stylization of 3D Gaussian Splatting（2026.06） 延续 3DGS 在"Multi-view stylization of Gaussian scenes"上的基本设定，主要改变是：基于半平衡最优传输的局部风格匹配与容量控制。3DGS 的显式高斯表示和可微渲染能力保持为基础。其设计动机是：把局部风格匹配重写为 semi-balanced optimal transport，并加入可调列容量约束以稳定风格分配。具体而言，针对每个渲染视图独立 2D feature matching 导致的视角不一致问题，在 3DGS 优化或编辑过程中约束风格特征的跨视角分配容量。适用边界是：最优传输约束可能增加优化成本；极端风格或复杂遮挡场景的局部一致性仍需关注。
\textbf{技术实现补充：} Capacity-Controlled Multi-View Stylization of 3D Gaussian Splatting。Capacity-Controlled Multi-View Stylization of 3D Gaussian Splatting（2026.06） 延续 3DGS 在"Multi-view stylization of Gaussian scenes"上的基本设定，主要改变是：基于半平衡最优传输的局部风格匹配与容量控制。3DGS 的显式高斯表示和可微渲染能力保持为基础。其设计动机是：把局部风格匹配重写为 semi-balanced optimal transport，并加入可调列容量约束以稳定风格分配。具体而言，针对每个渲染视图独立 2D feature matching 导致的视角不一致问题，在 3DGS 优化或编辑过程中约束风格特征的跨视角分配容量。适用边界是：最优传输约束可能增加优化成本；极端风格或复杂遮挡场景的局部一致性仍需关注。
\textbf{方法信息：} Tencent 参与了这篇论文。
\textbf{设计动机与边界：} 相对原始 3DGS 以独立高斯和图像重建误差为主的设定，Capacity-Controlled Multi-View Stylization of 3D Gaussian Splatting 把问题落在以下可实现的改动上：针对每个渲染视图独立 2D feature matching 导致的视角不一致问题，在 3DGS 优化或编辑过程中约束风格特征的跨视角分配容量。其设计目的不是单独增加一个模块，而是用上述表示、监督或渲染约束消除原流程在目标任务中的具体失配。把局部风格匹配重写为 semi-balanced optimal transport，并加入可调列容量约束以稳定风格分配。适用时还应保留论文给出的边界：最优传输约束可能增加优化成本；极端风格或复杂遮挡场景的局部一致性仍需关注。

\section{文本与语义引导编辑}

文本编辑路线的锚点是 InstructPix2Pix、SDS 和多视角 diffusion editing。3DGS 使局部高斯可被直接更新，但文本只给语义目标，不告诉系统哪些高斯该变、怎么在多个视角一致变。因此这些方法通常分三步：先定位编辑区域或语义对象，再生成多视角一致的 2D 编辑监督，最后把监督反传或重建回高斯。
% METHOD: 2311.14521
\paragraph{GaussianEditor（语义追踪与层次化喷洒）（2023.11）}
GaussianEditor~\cite{2311_14521} 是早期 3DGS 文本编辑锚点。它看到两个关键问题：第一，编辑目标在优化过程中会移动、分裂和致密化，静态 2D 掩码无法长期追踪；第二，SDS 等扩散指导梯度随机性大，直接更新高斯会不稳定。

方法为每个高斯增加语义掩码属性，初始由 SAM/DEVA 等 2D 分割结果监督，训练中持续维护语义标签；densification 时，新高斯继承父高斯标签，使目标对象在增殖后仍可被追踪。为了稳定扩散编辑，GaussianEditor 提出 Hierarchical Gaussian Splatting：不同层级控制不同粒度的高斯更新，低层级处理整体形状，高层级处理细节纹理。编辑策略包括文本引导 SDS、物体移除和区域集成。相对 3DitScene，它更像在已有 3DGS 上附加语义追踪和层次优化，而不是从单图构建语言解耦场景。
\textbf{实现细节：} GaussianEditor（2023.11） 是首个系统性 3DGS 编辑框架。其设计动机源于一个核心观察：3DGS 作为点云式显式表示，在面对扩散模型的随机梯度时缺乏神经网络的记忆与缓冲能力，容易导致高斯点产生剧烈而不稳定的更新。为解决这一问题，GaussianEditor 提出两个关键技术。其一是高斯语义追踪（Gaussian Semantic Tracing）：为每个高斯基元引入语义属性 $m_{ij}$（第 $i$ 个高斯对第 $j$ 个语义标签的隶属度），通过逆渲染将 2D 分割掩码反投影至 3D 高斯空间，权重计算为 $w_i^j = o_i(p) \cdot T_i^j(p) \cdot M^j(p)$，其中 $o_i$ 为不透明度、$T_i$ 为透射率、$M^j$ 为像素语义掩码。在自适应致密化过程中，新分裂的高斯继承父节点的语义标签，保证语义掩码在整个训练过程中持续有效——这克服了静态 2D/3D 掩码随场景演化而失效的问题。语义追踪还带来一个重要优势：在复杂场景中，2D 扩散模型往往难以有效编辑小物体；借助语义追踪，可以仅渲染目标物体并输入扩散模型，从而获得更精确的编辑监督。其二是层次化高斯泼溅（HGS, Hierarchical Gaussian Splatting）：按高斯的致密化轮次划分代际，初始高斯为第 0 代，第 $k$ 轮致密化产生的新高斯标记为第 $k$ 代。每代高斯施加不同强度的锚点损失以控制自由度： 其中 $P$ 为高斯的某一属性（位置 $x$、缩放 $s$、旋转四元数 $q$、不透明度 $\alpha$、颜色 $c$），$\hat{P}$ 为锚点状态，$\lambda_i$ 随代际递增——老一代高斯的锚点约束更强、自由度更低。直觉上，这一机制使早期代际逐渐"凝固"为稳定几何，而最新致密化的高斯在几乎无约束状态下雕琢细节，有效抑制了随机梯度导致的发散。
\textbf{锚点细化：} GaussianEditor：是率先提出的系统性 3DGS 编辑框架之一，针对扩散随机梯度使高斯产生剧烈不稳定更新的局限，用\textbf{高斯语义追踪 + 层次化高斯泼溅（HGS）}替代无约束迭代，核心思想是\textbf{按致密化轮次划分高斯代际，老一代锚点约束更强、最新代几乎无约束、雕琢细节}。直觉上，这一机制类似建筑中的``先打地基再砌墙''——早期代际承担整体几何骨架必须锁死防止漂移，最新代负责雕琢细节可以自由调整；语义追踪将 2D 分割掩码反投影至 3D 高斯空间、并在致密化时让新高斯继承父节点语义标签；HGS 对位置、缩放、旋转、不透明度、颜色五种属性分别施加代际递增的锚点损失，让早期代际逐渐``凝固''为稳定几何、而最新代精修细节。
\textbf{方法信息：} 代码：\url{https://github.com/buaacyw/GaussianEditor}；GitHub 关注度很高，约 1.4k stars；引用很多，Semantic Scholar 约 422 次。



% METHOD: 2404.18929
\paragraph{DGE（直接高斯编辑与极线一致图像编辑）（2024.04）}
DGE~\cite{2404_18929} 把 3D 编辑解耦为两个阶段：先得到多视角一致的 2D 编辑图，再直接优化 GS。它认为 SDS 或迭代式 diffusion update 的不稳定来自“边生成边优化”，因此先把图像编辑做干净。

第一阶段，从 3DGS 渲染一组沿相机轨迹排序的视图，用 InstructPix2Pix 编辑每个视图。为了让不同视图中的编辑区域对应同一 3D 部位，DGE 引入训练免费的极线一致约束：在一个视图的编辑区域，沿其他视图的极线搜索对应区域，约束编辑内容几何一致。第二阶段，用这些编辑后的多视图图像作为监督直接优化高斯参数。由于 3DGS 显式，方法还能融合多个视图的 2D 分割结果，识别哪些高斯需要局部编辑。相对 GSEdit，它更少依赖逐步 SDS，更多依赖几何约束下的一次性多视图编辑。
\textbf{实现细节：} DGE（Direct Gaussian Editor）（2024.04） 正是针对这一瓶颈提出了替代思路：与其在 3D 空间反复修补不一致的编辑结果，不如直接在图像空间保证多视角一致性，然后仅需将 3DGS 拟合到编辑后的图像即可完成 3D 编辑。与 GaussianEditor 通过 HGS 逐代凝固高斯属性来抑制迭代发散不同，DGE 将核心难题前移至 2D 编辑阶段——借鉴视频编辑中无需额外训练即可扩展图像生成器的思路，将多视角序列视为轨道视频，使多视角一致性等价于视频编辑中的时序一致性。DGE 将流程分为两步：关键视角编辑和特征注入。关键视角编辑阶段，在 IP2P 的自注意力层中引入时空注意力，使 $T$ 个视角的 query 同时关注所有视角的 key-value： 但该操作计算开销随视角数二次增长，因此实际执行中仅对随机采样的关键视角子集 $\mathcal{K} \subset \{1, \ldots, T\}$ 执行。在特征注入阶段，DGE 利用极线约束（Epipolar Constraint）将关键视角的编辑特征传播到其余视角：给定非关键视角 $t'$ 的特征图 $\mathbf{F}_{t'}$ 和最近关键视角 $k^{*}$ 的特征图 $\mathbf{F}_{k^{*}}$，对应关系通过基础矩阵 $F$ 约束搜索空间获得： 其中 $D$ 为余弦距离，$Fv \parallel u$ 表示点 $v$ 在关键视角中的对应点必须位于极线 $Fu$ 上。这一约束将对应搜索从二维平面压缩到一维极线上，降低了误匹配率。最终 DGE 通过渲染损失 $\theta' = \arg\min_\theta \sum_{t=1}^{T} \| I'_t - R(\theta, \pi_t) \|$ 直接优化 3DGS 参数，整个编辑流程无需额外训练扩散模型。
\textbf{锚点细化：} DGE：针对 GaussianEditor 仍在 3D 空间修补不一致编辑、IP2P 单目编辑器跨视角独立采样的局限，用\textbf{图像空间多视角一致编辑}替代 3D 反复修补，核心思想是\textbf{将多视角序列视为轨道视频，使多视角一致性等价于视频编辑的时序一致性}——直觉上，绕一个物体拍的多视角序列本质就是一段``绕轨道飞行的视频''，视频编辑领域已有时序注意力成熟工具保证帧间一致；关键视角编辑阶段在 IP2P 自注意力层引入时空注意力让 $T$ 个视角 query 同时关注所有视角 key-value，特征注入阶段利用极线约束（同一 3D 点在不同视角的像素必在对极线上）将关键视角编辑特征沿一维极线传播到其余视角，将一致性难题前移至 2D 编辑阶段。
\textbf{方法信息：} 代码：\url{https://github.com/silent-chen/DGE}；GitHub 关注度较高，约 129 stars；引用很多，Semantic Scholar 约 104 次。



% METHOD: 2405.18424
\paragraph{3DitScene（语言引导的对象级解耦高斯）（2024.05）}
3DitScene~\cite{2405_18424} 是文本/语义编辑路线中值得详细读的锚点。它处理的是单张图像生成出的 3D 场景如何被语言编辑的问题：如果只用 InstructPix2Pix 或 SDS 在渲染图上改，系统知道“想改成什么”，却不知道三维场景里“哪一组高斯对应这个对象”。3DitScene 的改变是先把语言语义写进高斯表示，使对象可以被文本查询选中，再对选中的高斯组做移动、删除、形变或外观修改。

完整流程可以分成四步。第一步是单图到 3DGS 初始化：用深度预测把输入图像提升为点云，并以此初始化高斯，而不是从随机场开始优化；随后用重建损失保持输入视角内容，用 Stable Diffusion 的 SDS 约束补足外观合理性。第二步是视角扩展：当新相机看到原图中没有覆盖的区域时，先用扩散修复生成新内容，再用深度预测把新像素提升到三维；论文中特别处理了新旧深度对齐，避免简单拼接造成深度断层和场景变形。第三步是语义蒸馏：系统渲染当前 3DGS 视图，在 2D 图像上提取 CLIP/OpenCLIP 等语言特征，再通过可微 splatting 把这些语义监督分配回可见高斯，使每个高斯不仅有颜色和形状，还有可查询的语义特征。第四步才是编辑：文本查询先选出目标对象的高斯集合，之后可以对该集合执行移动、缩放、删除、外观替换，或对整场景做风格/光照变化。

因此，3DitScene 相对 Instruct-NeRF2NeRF 的增量不只是把编辑对象换成 3DGS，而是把“渲染图编辑”改成“语义高斯集合编辑”；相对 GaussianEditor 这类依赖编辑时掩码传播的方法，它把对象定位前置为表示学习的一部分。读者后面看语义高斯编辑、开放词汇对象操控和单图场景编辑时，可以把 3DitScene 当作一条基线：先构建可语言寻址的显式 3D 场，再用高斯集合承接编辑指令。
\textbf{实现细节：} 语言驱动解耦：3DitScene。3DitScene（2024.05） 保留了 Point'n Move 的对象级操控能力，但将分割方式从"用户点击 → SAM 分割"替换为"语言/框选 → CLIP 语义匹配"，使用户无需精确点击即可选中目标。具体地，3DitScene 为每个 3D 高斯基元额外嵌入一个 CLIP 语言向量 $e \in \mathbb{R}^C$，利用 3DGS 的可微分泼溅将该向量渲染为 2D 语义特征图 $E \in \mathbb{R}^{C \times H \times W}$，以蒸馏损失。
\textbf{实现补充：} 实现上还要核对输入、输出、参数、特征、初始化、训练损失、正则、渲染或推理接口，以及优化中的剪枝、分裂、采样、编码或解码步骤；这些抓手决定该方法到底改了 3DGS 管线的哪一环。
\textbf{方法信息：} ByteDance 参与了这篇论文；代码：\url{https://github.com/zqh0253/3ditscene}；GitHub 关注度较高，约 259 stars。



% METHOD: 2412.00518
\paragraph{Instant3dit（多视角修补式快速编辑）（2024.12）}
Instant3dit~\cite{2412_00518} 把 3D 编辑重新表述为多视角图像 inpainting。SDS 类方法每次编辑都要优化 3D 表示，慢且容易模糊；Instant3dit 的改变是训练一个专门的多视角扩散修补模型，在约数秒内给出一致的编辑视图，再重建回 3D。

输入是 3D 形状或高斯场景、3D 掩码区域和文本描述。系统从多个视角渲染形状和掩码，组成 $2\times2$ 多视角网格图像，交给微调后的扩散模型一次性修补。微调时用多种掩码模拟用户编辑需求，使模型学会在多个视图间保持一致；推理时不再对每个编辑目标做 SDS 优化，而是先得到一致的多视角编辑图，再用重建模型或少量优化把结果写回 3D 表示。修补完成后，可用 GaussianSRM、NeRF-LRM 或优化方法把多视角结果转回 3DGS 或其他 3D 表示。相对 GaussianEditor 的逐步 SDS，它把生成难题放到多视角 2D 模型中解决；相对 DGE，它依赖训练过的多视角修补先验。
\textbf{设计动机与边界：} 相对原始 3DGS 以独立高斯和图像重建误差为主的设定，Instant3dit 把问题落在以下可实现的改动上：Instant3dit的核心洞察是将3D编辑重新表述为多视角图像修补问题。输入包括3D形状S、3D掩码区域M和文本描述y。首先从多个视角渲染3D形状和掩码，生成2x2网格的多视角图像。然后利用经过特殊微调的扩散模型进行视角一致的修补。微调策略包括多种掩码设计以模拟用户的编辑操作。修补完成后，利用大型重建模型（如NeRF-LRM、GaussianSRM或优化方法）将多视角结果重建回3D表示。该方法避免使用SDS类优化，从根本上解决了SDS导致的模糊和低质量问题。其设计目的不是单独增加一个模块，而是用上述表示、监督或渲染约束消除原流程在目标任务中的具体失配。提出将3D编辑转化为多视角图像修补问题，利用微调的扩散模型在约3秒内完成编辑，无需SDS优化，质量显著优于现有方法，支持网格、NeRF和高斯Splat等多种表示。适用时还应保留论文给出的边界：编辑效果受限于扩散模型的生成能力；复杂几何编辑（如拓扑变化）可能不够精确；掩码设计需要一定用户交互。
\textbf{方法信息：} Adobe 参与了这篇论文；引用很多，Semantic Scholar 约 37 次。



% METHOD: 2412.08152
\paragraph{ProGDF（一次训练后的实时可控编辑）（2024.12）}
ProGDF~\cite{2412_08152} 的问题是编辑交互速度。传统 3DGS 编辑是 in-the-loop：用户改一次，系统重新跑一段优化；ProGDF 把它改成 out-of-loop：先训练一个能预测编辑差分的场，之后用户用滑动条实时调编辑强度。

方法先用 Progressive Gaussian Splatting 记录编辑优化轨迹，在多个时间步保存高斯状态快照，并用渲染质量正则保证中间结果可用。然后训练 Gaussian Differential Field，一个轻量网络以空间位置和时间/强度控制量为输入，输出该位置高斯属性的偏移。训练完成后，界面上的控制信号会即时查询 GDF 并更新高斯，无需再次优化，单次调整可到 $0.02$ 秒量级。相对 DGE 或 GaussianEditor，它牺牲了一些任意新编辑的开放性，换来同一编辑族中的实时可控。
\textbf{实现细节：} 近期工作从架构层面重新思考 3DGS 编辑范式。前三条路线（SDS 引导、迭代数据更新、渐进式优化）共享同一个范式瓶颈：编辑过程与优化过程交织在一起，即 in-the-loop 训练——每次编辑指令的修改都需要重新启动一轮完整优化（据 ProGDF 原文每次调整需数十分钟训练成本），用户无法实时交互。ProGDF（2024.12） 从根本上扭转了这一范式，将编辑从 in-the-loop 转变为 out-of-loop：训练一次、实时编辑。动机。in-the-loop 训练的编辑范式使用户每次修改指令都需等待完整优化，无法实现交互式体验；ProGDF 用轻量级神经网络学会从原始场景到目标场景的整个编辑轨迹，用户只需拖动滑块选择轨迹上的时间点即可实时获得编辑结果。其直觉是：如果用一个轻量级神经网络学会从原始场景到目标场景的整个编辑轨迹，那么用户只需拖动滑块选择轨迹上的某个时间点，就能在 $0.02$ 秒内获得对应的编辑结果，无需再次优化。ProGDF 由两个组件协同实现这一目标：渐进式高斯泼溅（PGS, Progressive Gaussian Splatting）负责在优化轨迹上提取高质量中间结果；高斯微分场（GDF, Gaussian Differential Field）负责学习这些中间结果并参数化为滑块控制信号。PGS 的核心挑战在于：常规优化轨迹的前期大量高斯朝不精确方向更新，导致纹理模糊和细节扭曲；后期虽然收敛但中间状态缺乏多样性，难以提供丰富的训练监督。为此，PGS 提出两项约束平滑优化轨迹。渐进约束 $L_{prog} = \beta^{t/s}\sum_{i=1}^{N}\|\delta_i^t\|$ 通过衰减系数 $\beta^{t/s}$（$0<\beta<1$，$s$ 用于缓解指数爆炸）自适应压制每步更新幅度，使场景从原始到目标的过渡更加平稳。渲染质量正则化 $L_{render} = -\|\nabla_{Laplacian}^{}2 R(\Theta^t,v)\|^{}2$ 利用拉普拉斯梯度衡量图像清晰度，梯度越大表示越清晰，以此显式约束中间结果的纹理质量。
\textbf{方法信息：} Tencent 参与了这篇论文。
\textbf{设计动机与边界：} 相对原始 3DGS 以独立高斯和图像重建误差为主的设定，ProGDF 把问题落在以下可实现的改动上：ProGDF包含两个关键组件：Progressive Gaussian Splatting（PGS）和Gaussian Differential Field（GDF）。PGS引入渐进约束提取编辑过程的多样化中间结果，并使用渲染质量正则化提高中间结果质量。具体地，在3DGS优化轨迹上均匀采样多个时间步，保存高斯状态的快照作为中间结果。GDF是一个轻量神经网络，以空间位置为输入，估计该位置处因编辑导致的高斯属性偏移。GDF利用PGS提供的多时间步中间结果作为监督进行训练，同时将时间步作为控制信号赋予编辑可控性。训练完成后，用户可通过滑动条等界面实时调整编辑强度和效果，每次调整仅需0.02秒。其设计目的不是单独增加一个模块，而是用上述表示、监督或渲染约束消除原流程在目标任务中的具体失配。提出Progressive Gaussian Differential Field（ProGDF），将3D编辑从传统的in-the-loop训练范式转变为out-of-loop训练。仅需一次训练，即可在0.02秒内通过用户界面实时调整编辑结果，实现可控3D编辑和灵活细粒度操作，消除复杂超参数调整和重新训练的需要。适用时还应保留论文给出的边界：GDF的表达能力受限于网络容量；对极端编辑操作可能效果有限；编辑范围受限于原始训练过程中探索的编辑轨迹。
% METHOD: 2403.05154
\paragraph{GSEdit（IP2P 与 SDS 驱动的文本编辑）（2024.03）}
GSEdit~\cite{2403_05154} 是文本编辑的基础方法之一。输入可以是已有 3DGS，也可以是从网格多视角渲染或生成模型输出重建的高斯；用户给出文本指令，系统用 InstructPix2Pix 和 SDS 逐步优化高斯参数。

每次迭代从随机相机渲染当前 3DGS 图像，把图像和文本提示送入 IP2P 得到编辑目标或噪声预测，再把 SDS 梯度通过可微 rasterization 反传到高斯颜色、透明度和几何属性。为了避免编辑一步过猛破坏场景，它使用渐进式方向增强，在 CLIP/扩散空间中逐步推动图像朝目标语义变化。多视角一致性不是由显式匹配保证，而是通过不断从不同相机采样并更新同一组高斯自然形成。相对 DGE，它更灵活但更慢；相对 Instant3dit，它不需要专门训练多视角修补模型。
\textbf{实现细节：} 动机。3DGS 缺乏神经网络的记忆与缓冲能力，在面对扩散模型的随机梯度时容易产生不稳定更新，直接套用 NeRF 的 SDS 编辑范式效果有限。以 DreamFusion 的 SDS 为基线，其核心思想是：在每个训练步随机采样相机 $p$ 渲染图像 $I^p_{\mathrm{RGB}}$，对渲染结果加噪后用扩散模型的 U-Net 预测噪声残差，再将残差梯度反传至 3DGS 参数。GSEdit（2024.03） 将此框架适配到 3DGS 编辑场景：采用 InstructPix2Pix（IP2P）作为编辑扩散模型——IP2P 是一种条件扩散模型，同时接受原始图像和文本指令作为输入，能在保持图像整体布局的同时按指令修改局部内容。在每个训练步，GSEdit 随机采样相机渲染当前场景，将渲染图像与编辑文本提示一同输入 IP2P。IP2P 的噪声预测与实际添加噪声之间的残差被用来引导 3DGS 参数的更新方向。直觉上，这个残差告诉优化器"当前渲染结果距离'按文本编辑后的样子'还差多少"。其 SDS 损失定义为 其中 $z_t$ 是对渲染图像添加噪声 $\epsilon$ 后的隐编码（latent code），$\hat{\epsilon}_\phi$ 是 IP2P U-Net 在给定原始图像条件 $c_I$ 和编辑文本 $c_T$ 下预测的噪声，$w(t)$ 是时间步权重。残差项 $\hat{\epsilon}_\phi - \epsilon$ 编码了 IP2P 认为图像需要被编辑的程度与实际噪声之间的差异，该差异通过链式法则反传至 3DGS 的位置、颜色、透明度等参数。GSEdit 的一个关键设计是时间步截断策略：将 $t$ 限制在 $[t_{\min}, t_{\max}]$ 区间内，并随训练推进线性降低该区间（参考 DreamTime 的策略）。之所以需要截断，是因为全时间步采样中高 $t$ 值对应的大噪声会引入过度随机扰动，容易破坏已有的几何结构。截断后，初始阶段加噪量较大（典型设置 $t_{\min}=0.02$, $t_{\max}=0.98$），允许模型在粗粒度上调整整体结构和外观。随着训练推进，$t_{\max}$ 逐步降低，加噪量减小，优化焦点转向细粒度纹理精修。整个编辑过程仅需约 500 步迭代，在消费级 GPU 上几分钟内即可完成，比同期 NeRF 编辑方法（如 Instruct-NeRF2NeRF 的约一小时）快一个数量级。
\textbf{方法信息：} 引用很多，Semantic Scholar 约 43 次。
\textbf{设计动机与边界：} 相对原始 3DGS 以独立高斯和图像重建误差为主的设定，GSEdit 把问题落在以下可实现的改动上：GSEdit管线包含四步：(1)输入表示：支持任意高斯泼溅场景输入，包括从3D网格的多视角渲染重建的GS模型和DreamGaussian等生成模型的输出。(2)编辑过程：每步从随机相机渲染RGB图像，通过IP2P（InstructPix2Pix）模型结合编辑文本提示处理渲染图。采用单步调整策略在CLIP空间中进行渐进式方向增强，保持场景完整性。处理后通过可微渲染更新高斯参数，然后重新渲染更新图像。此循环过程持续到编辑收敛。(3)梯度反向传播：将SDS损失适配到GS编辑中，通过IP2P U-Net的噪声预测计算梯度，反向传播到高斯参数。(4)多视角一致性：通过循环渲染-编辑-更新过程自然保持多视角一致性。其设计目的不是单独增加一个模块，而是用上述表示、监督或渲染约束消除原流程在目标任务中的具体失配。GSEdit提出基于3DGS的文本引导3D物体编辑管线，利用预训练图像扩散模型（IP2P）逐步优化高斯参数，通过SDS损失确保编辑精确性，支持从网格或高斯生成模型输入，在消费级硬件上几分钟内完成编辑，显著快于NeRF编辑方法。适用时还应保留论文给出的边界：主要适用于前景物体编辑；几何变化编辑能力有限（偏向纹理修改）；编辑质量受IP2P模型能力限制；可能需要多次尝试获得理想编辑效果
% METHOD: 2507.05426
\paragraph{Mastering Regional 3DGS（区域定位、初始化与序贯编辑）（2025.07）}
Mastering Regional 3DGS~\cite{2507_05426} 专门处理局部区域编辑。很多文本编辑方法的问题不是扩散模型不会改图，而是不知道 3D 场景中哪一块该新增、替换或细化。该方法把区域编辑拆成定位、初始化和细化三步。

定位阶段使用 InstructPix2Pix，在条件化和非条件化噪声预测之间取差异，得到每个视图的 2D 编辑区域，再反渲染到 3D 空间，标记需要修改的高斯。初始化阶段选择与文本最相关、掩码面积最大的首选视图，用 Depth Anything 预测深度，并用未编辑高斯渲染深度校准尺度，在目标区域内生成粗 3DGS 点。细化阶段按条件序贯视角逐步加入更多视图，每一步以上一步结果为条件继续编辑和优化，逐渐补结构和纹理。相对 GaussianEditor，它更强调新增或局部替换时的几何初始化。
\textbf{实现细节：} 沿渐进式定位路线，后续方法从不同角度持续改进。Mastering Regional 方法（2025.07） 利用 2D 扩散编辑精确定位 3DGS 修改区域，结合深度先验初始化和序贯视角细化，在 SOTA 性能上实现 4 倍加速。互补与一致性信息方法（2025.03） 利用多视角互补信息增强编辑一致性。
\textbf{方法信息：} Tencent 参与了这篇论文。
\textbf{设计动机与边界：} 相对原始 3DGS 以独立高斯和图像重建误差为主的设定，Mastering Regional 3DGS 把问题落在以下可实现的改动上：方法包含三个核心步骤：(1) 2D到3D定位：使用IP2P(InstructPix2Pix)计算UNet在条件化（图像+文本）和非条件化（仅图像）下的噪声预测差异，定位每视角的2D编辑区域，然后反渲染到3D空间获得稀疏3D掩码标记需修改的高斯点；(2) 基于预测深度的增量初始化：选择与文本最相关的视角（最大掩码面积）作为首选视角，使用Depth Anything模型预测深度，通过未编辑高斯点的渲染深度校准尺度偏差，在识别的3D区域内用校准后的深度初始化粗3DGS点；(3) 条件序贯视角细化：从首选视角开始，逐步增加新视角进行迭代编辑和优化，每步条件化于之前视角的编辑结果，逐渐增强结构细节和纹理，确保跨视角一致性。其设计目的不是单独增加一个模块，而是用上述表示、监督或渲染约束消除原流程在目标任务中的具体失配。提出基于2D扩散编辑的区域3DGS编辑方法，通过2D到3D反渲染定位、深度先验初始化粗3DGS、条件序贯视角细化迭代优化，实现SOTA性能同时4倍加速。适用时还应保留论文给出的边界：依赖2D扩散模型和深度估计模型的精度；对大范围结构变化的编辑能力有限；深度校准在校准区域外可能不准确。
% METHOD: 2508.16932
\paragraph{Invert3D（3D 表示到文本嵌入空间的对齐）（2025.08）}
Align 3D Representation and Text Embedding for 3D Content Personalization~\cite{2508_16932} 可理解为 3D 内容的 textual inversion。普通文本编辑把文本提示作用到渲染图上，却没有一个代表“这个 3D 对象本身”的文本嵌入；Invert3D 的改变是学习 camera-conditioned text inversion，把 3DGS 或 NeRF 场景映射到可编辑的文本嵌入空间。

具体做法从多个视角渲染 3D 场景，用 VAE 编码器得到潜空间表示，并优化一个 3D 嵌入 $z^*$，使其在不同相机条件下都能解释对应渲染图。相机参数作为条件输入，帮助嵌入保持 3D 多视角一致，而不是退化为某一张图的 2D token。得到嵌入后，用户可以用自然语言提示在潜空间中做个性化编辑，再通过 3DGS 渲染输出。相对 3DitScene 的高斯级语义蒸馏，Invert3D 更偏把整个 3D 内容注册进文本生成模型的嵌入空间。
\textbf{实现细节：} 在个性化编辑方向，CP-GS（2025.05） 通过预训练 image-to-3D 模型和迭代 LoRA 微调实现单图参考的 3DGS 个性化。Invert3D（2025.08） 首次建立 3DGS 与文本嵌入空间的对齐，通过相机条件化逆映射实现自然语言引导的 3D 内容个性化编辑 。这些方法共同展示了文本引导编辑从"修改现有场景"向"个性化定制场景"的演进趋势。
\textbf{技术实现补充：} 在个性化编辑方向，CP-GS（2025.05） 通过预训练 image-to-3D 模型和迭代 LoRA 微调实现单图参考的 3DGS 个性化。Invert3D（2025.08） 首次建立 3DGS 与文本嵌入空间的对齐，通过相机条件化逆映射实现自然语言引导的 3D 内容个性化编辑 。这些方法共同展示了文本引导编辑从"修改现有场景"向"个性化定制场景"的演进趋势。
\textbf{方法信息：} NVIDIA 参与了这篇论文。
\textbf{设计动机与边界：} 相对原始 3DGS 以独立高斯和图像重建误差为主的设定，Align 3D Representation and Text Embedding for 3D Content Personalization 把问题落在以下可实现的改动上：方法核心是相机条件化文本逆映射(CTI)：给定3D场景（3DGS或NeRF表示），通过可微渲染从多个视角渲染图像，使用VAE编码器编码到潜空间，然后优化一个3D嵌入z*使其与各视角渲染结果在潜空间中一致。相机参数作为显式输入确保嵌入保持3D结构的多视角一致性。优化目标为最小化渲染图像编码与文本条件解码的MSE距离。获得3D嵌入后，可通过自然语言提示在潜空间中直接操作实现个性化编辑，再通过3DGS渲染输出结果。其设计目的不是单独增加一个模块，而是用上述表示、监督或渲染约束消除原流程在目标任务中的具体失配。提出Invert3D框架，首次建立3D表示与文本嵌入空间的对齐。设计相机条件化3D到文本逆映射机制(Camera-conditioned Text Inversion, CTI)，将3D场景（3DGS或NeRF）投影到与文本嵌入对齐的3D嵌入空间中，实现通过自然语言提示直接个性化和编辑3D内容，无需昂贵的知识蒸馏重训练过程。适用时还应保留论文给出的边界：3D嵌入的表达能力可能不足以捕捉极其复杂的场景细节；编辑效果的精细控制仍有限；对训练时未见过的编辑类型泛化能力待验证。

\section{交互式物体操控}

交互式编辑的锚点不是文本，而是用户直接给出的 2D/3D 操作：拖拽点、画掩码、涂色或选择部件。3DGS 的显式点状结构很适合响应这些操作，但高斯不是 mesh 顶点，没有天然拓扑；因此方法要么借 diffusion 修正变形后的内容，要么借部件或体素特征把局部编辑传播到相邻高斯。
% METHOD: 2501.18672
\paragraph{Drag Your Gaussian（拖拽式几何编辑）（2025.01）}
Drag Your Gaussian~\cite{2501_18672} 是交互式几何编辑的锚点之一。文本编辑适合说“把它变成某种风格”，但很难精确表达“把这个把手拖到那个位置”或“让这条腿向左弯”。DYG 的输入因此不是一句文本，而是已有 3DGS、用户圈出的 3D 掩码，以及若干 handle-target 控制点对；输出是在局部区域完成拖拽后的高斯场景。

方法先解决交互输入问题。用户在 GUI 中从多个视角选择待编辑区域，系统把这些视锥约束合成为三维掩码；用户再给出若干 3D handle-target 点对，明确局部几何应当向哪里移动。直接移动被选高斯会遇到一个问题：3DGS 是稀疏点状表示，目标区域可能没有足够高斯承接连续形变。DYG 因此加入 multi-scale triplane positional encoder 和 region-specific positional decoder。三平面可以看作轻量的连续空间支架：任意三维位置投影到 $xy$、$xz$、$yz$ 三个平面查询多尺度特征，再由解码器预测掩码内高斯的位移；这样即使目标区域原本高斯稀疏，位移场也能连续传播。

第二个关键是把 2D 拖拽扩散先验接到 3D 优化里。DYG 保留初始高斯的镜像副本，从采样相机渲染 RGB 和掩码，把 3D 控制点投影成 2D 拖拽条件，再用拖拽式 latent diffusion 模型提供图像层面的编辑方向。Drag-SDS 把这个 2D 条件先验变成可反传的 score distillation 损失，与三维控制点约束一起优化目标区域；Soft Local Edit 约束掩码外内容，减少局部编辑向全场景扩散。相对 GSDeformer 的 cage/代理结构，DYG 不要求预先构建网格笼；相对后续 3DGS-Drag，它更强调三平面连续支架和 Drag-SDS 的组合。论文展示了变形、平移和形状改变等场景，实验还包含 Mip-NeRF360、Instruct-NeRF2NeRF 场景与用户研究。
\textbf{锚点细化：} DYG：针对笼式与网格代理方法均依赖额外代理结构、直接在 3DGS 上拖拽时目标区域往往高斯基元稀疏导致拖拽后无法稳定锚定的局限，用\textbf{多分辨率三平面位置编码器作为连续隐式几何支架}替代显式高斯直接拖拽，核心思想是\textbf{将 3D 空间投影到三个正交平面 $(xy,xz,yz)$ 上、在多个分辨率上维护可学习特征图，任意位置 $p$ 的特征为三平面多尺度拼接 $\mathbf{f}(p)=\phi(\bigodot_s\bigodot_c\psi^s(H_c,\pi_c(x,y,z)))$}。直觉上，``三平面''是显式 3D 体素的特征降维版——存一个 $N^3$ 的体素特征图代价太高，而把它分解为三张 $N^2$ 的特征图（xy、xz、yz 三个视角的``切片堆叠''）只需 $3N^2$ 存储，再通过双线性插值就能查询空间中任意位置的特征；这一连续特征场为目标区域空白位置提供合理编码，使高斯可在任意位置获得可靠位移预测。区域感知位置解码器区分掩码内/外高斯分别预测位移；Drag-SDS 损失将用户拖拽意图编码为 2D 拖拽扩散模型条件输入。
\textbf{方法信息：} Tencent、Baidu 参与了这篇论文；代码：\url{https://github.com/quyans/drag-your-gaussian}；GitHub 关注度较高，约 165 stars。



% METHOD: 2512.05354
\paragraph{SplatPainter（带测试时训练的交互式高斯绘制）（2025.12）}
SplatPainter~\cite{2512_05354} 面向 3D 高斯资产的交互式精修。它的锚点是 Gaussian LRM 前馈重建，但普通前馈模型输出的高斯冗余且缺少可编辑潜表示；SplatPainter 的改变是先把高斯压缩成紧凑的体素化特征表示，再用带 Test-Time Training 的精修模块响应用户输入。

具体流程先从多视角渲染图生成带丰富特征的高斯云，再把高斯体素化为局部单元，每个体素包含 token 和高斯属性。局部体素 Transformer 先在层内处理，再跨层聚合，得到上下文感知且紧凑的 latent。编辑时，前馈模型直接预测高斯属性更新，包括位置、尺度、不透明度、旋转和 SH；TTT 层把用户新输入融入当前状态，使系统能多轮适应。它支持局部细节精修、涂鸦式局部绘画和全局重着色。相对 Instant Colorization 的线性反投影，SplatPainter 更偏学习式资产编辑器。
\textbf{实现细节：} 沿多视角扩散修复路线，后续方法将能力扩展到更复杂的场景类型。AuraFusion360（2025.02） 是首个针对 360 度无界场景的参考式 3D 修复方法，通过深度感知未见掩码生成、自适应引导深度扩散和 SDEdit 细节增强实现高质量对象移除。LightHarmony3D（2026.03） 专门解决物体插入中的光照一致性问题，通过生成式 HDR 环境图预测和 PBR 引导的阴影合成确保物理一致的阴影整合。
\textbf{实现补充：} 进一步拆开来看，输入侧对应图像、点云、位姿、语义提示或已训练高斯资产；输出侧对应新的高斯参数、渲染图像、编辑结果、地图状态或评测指标。表示层需要说明位置、尺度、旋转、不透明度、特征、材质或 MLP 解码器怎样改；训练层需要说明损失、正则、采样、剪枝、分裂或压缩何时发生；推理层则看渲染、查询、规划或系统接口如何接到原始 3DGS / Scaffold-GS / 4DGS 管线上。
\textbf{方法信息：} Adobe 参与了这篇论文。
% METHOD: 2601.07963
\paragraph{3DGS-Drag（显式拖拽与扩散校正）（2026.01）}
3DGS-Drag~\cite{2601_07963} 同样处理拖拽编辑，但它更强调在显式高斯层面直接给出几何引导。输入是预训练 3DGS 和若干控制点到目标点的对应关系；系统先用复制、移动和变形高斯的方式生成多视角一致的粗几何，再让扩散模型修正图像伪影。

具体流程包含变形引导、扩散引导和渐进式编辑。变形引导利用高斯中心和显式参数，把被选中区域按照控制点位移进行局部形变或内容扩展；扩散引导在渲染视图上修补因拉伸、空洞或遮挡产生的不自然内容；渐进策略把大幅编辑拆成多个小步，并在每一步重新渲染视图、更新拖拽配置和高斯参数。公开实现的入口以 \texttt{train\_drag.py} 和配置文件组织，说明它仍是“拖拽条件下的场景优化”，不是一次前馈编辑器。相对 Drag Your Gaussian，3DGS-Drag 的三维操作更直接，代价是通常需要更多逐步优化时间。
\textbf{方法信息：} 代码：\url{https://github.com/dongjiahua/3dgs-drag}。
\textbf{设计动机与边界：} 相对原始 3DGS 以独立高斯和图像重建误差为主的设定，3DGS-Drag 把问题落在以下可实现的改动上：3DGS-Drag框架以预训练的3DGS模型和若干控制点-目标点对为输入。核心包括：(1)变形引导：利用3DGS的显式表示，以复制粘贴方式直接变形3D高斯，为多视角提供一致的几何引导；(2)扩散引导：使用扩散模型对变形后渲染图像进行内容校正，移除伪影并生成正确内容；(3)渐进式编辑策略：将大幅编辑分解为多步渐进过程，用户可在中间步骤停止。该方法无需先验信息，通过两重引导从不同视角显式正则化编辑一致性。其设计目的不是单独增加一个模块，而是用上述表示、监督或渲染约束消除原流程在目标任务中的具体失配。提出3DGS-Drag，首个弥合变形方法和2D编辑方法之间差距的3D拖拽编辑框架。利用3DGS的显式表示进行变形引导实现一致的几何修改，结合扩散引导进行内容校正和视觉质量增强，支持包括运动变化、形状调整、修复和内容扩展在内的广泛编辑操作，在单块RTX 4090上10-20分钟完成编辑。适用时还应保留论文给出的边界：编辑时间仍较长（10-20分钟）；大幅编辑可能需要多步渐进操作；对非常复杂的拓扑变化处理可能受限。
% METHOD: 2604.17155
\paragraph{Instant Colorization of Gaussian Splats（正规方程反投影着色）（2026.04）}
Instant Colorization of Gaussian Splats~\cite{2604_17155} 解决的是如何把 2D 信息快速写回已有高斯场。常规做法是把目标颜色、特征或掩码渲染成监督，再用梯度下降优化高斯；该方法把它写成可见性加权最小二乘问题，用正规方程直接求解。

它先用 3DGS renderer 计算每个高斯对每个像素的可见性权重，这个权重自然包含遮挡、透明度和 alpha blending 贡献。然后对 RGB、DINO 特征或分割掩码等通道分别建立线性系统，求解哪些高斯属性最能解释 2D 输入。由于不需要迭代反传，它比梯度下降快一个量级，适合交互式重光照、特征丰富化和 SAM2 掩码到 3D 分割。相对 ReCoGS，它更通用：不是只做区域重着色，而是把任意 2D 通道映射到高斯。
\textbf{实现细节：} 交互式操控还有多种拓展方向。在参考图引导方面，ICE-G（2024.06） 通过 SAM 分割和 DINO（自监督视觉 Transformer）语义匹配实现单张参考图像编辑。MiraGe（2024.10） 利用镜面反射将 2D 图像映射到 3D 空间。在特定创作形式方面，SandSim（2026.04） 通过曲线引导高斯表示从单张图像重建沙画创作过程。其核心做法是将 3DGS 的线性结构用曲线参数化，从单张图像中恢复沙画艺术家的创作轨迹。在效率与专用工具方面，Instant Colorization（2026.04） 用正规方程实现高效 2D 到 3D 信息映射（比梯度下降快一个数量级）。SketchFaceGS（2026.04） 是首个基于草图的实时 3D 高斯头部生成与编辑框架 。 %=====================================================================。
\textbf{技术实现补充：} 交互式操控还有多种拓展方向。在参考图引导方面，ICE-G（2024.06） 通过 SAM 分割和 DINO（自监督视觉 Transformer）语义匹配实现单张参考图像编辑。MiraGe（2024.10） 利用镜面反射将 2D 图像映射到 3D 空间。在特定创作形式方面，SandSim（2026.04） 通过曲线引导高斯表示从单张图像重建沙画创作过程。其核心做法是将 3DGS 的线性结构用曲线参数化，从单张图像中恢复沙画艺术家的创作轨迹。在效率与专用工具方面，Instant Colorization（2026.04） 用正规方程实现高效 2D 到 3D 信息映射（比梯度下降快一个数量级）。SketchFaceGS（2026.04） 是首个基于草图的实时 3D 高斯头部生成与编辑框架 。 %=====================================================================。
\textbf{方法信息：} 代码：\url{https://github.com/dlieber01/instant-colorization-of-gaussian-splats}。
\textbf{设计动机与边界：} 相对原始 3DGS 以独立高斯和图像重建误差为主的设定，Instant Colorization of Gaussian Splats 把问题落在以下可实现的改动上：算法核心：(1) 可见性计算：利用可微分渲染器计算每个高斯在每个像素的可见性，自动整合透明度和遮挡因素；(2) 正规方程求解：将着色问题建模为可见性加权最小二乘问题，通过正规方程直接求解而非迭代优化；(3) 多通道扩展：支持RGB颜色通道、神经特征和分割掩码等多种信息类型；(4) 三种应用场景：场景重光照（不同光照条件下重新着色）、特征丰富化（用DINOv2提取的特征增强3DGS）、3D语义分割（用SAM2掩码进行3D分割）。其设计目的不是单独增加一个模块，而是用上述表示、监督或渲染约束消除原流程在目标任务中的具体失配。提出一种高效算法将2D图像信息（颜色、神经特征、分割掩码）映射回已有3DGS场景。使用正规方程求解可见性加权最小二乘问题，显式处理遮挡、半透明混合和视角相关效果。比梯度下降方法快一个数量级，支持场景重光照、特征丰富化和3D语义分割。适用时还应保留论文给出的边界：依赖已有场景的质量；对极端视角相关效果处理有限；大规模场景的内存开销。

\section{几何变形与操控}

几何变形路线的共同问题是：3DGS 没有 mesh 的边和面，直接移动高斯容易造成局部断裂。解决思路通常是引入代理几何，例如 cage、triangular mesh、SDF surface 或粗 mesh，把用户的形变先作用在有拓扑的结构上，再传播给高斯的位置、旋转和尺度。
% METHOD: 2405.15491
\paragraph{GSDeformer（笼式变形）（2024.05）}
GSDeformer~\cite{2405_15491} 是 3DGS cage-based deformation 的基础方法。它不修改原始 3DGS 表示，也不要求重训练，而是把已有高斯中心当作代理点云，通过 cage 坐标传递形变。

具体流程先自动生成包围 3DGS 的三角 cage：从高斯渲染多视角图像，用 SfM 或重建方法得到网格，再通过带防翻转约束的 QEM 边折叠简化成低多边形 cage。接着从高斯中心生成代理点云，为每个代理点计算 cage 坐标。用户编辑 cage 顶点后，代理点根据坐标得到新位置，再把仿射变换传给对应高斯的位置、旋转和缩放。当形变引起强弯曲时，方法把单个高斯分裂为多个子高斯，让不同子高斯贴近弯曲后的局部形状。相对 Drag Your Gaussian，GSDeformer 不依赖扩散先验，强调实时、可解释和兼容已有模型。
\textbf{实现细节：} GSDeformer（2024.05） 是首个笼式 3DGS 变形方法，其核心思路是将经典均值坐标（Mean Value Coordinates, MVC）适配到离散高斯基元。直觉上，均值坐标为笼内部的每个点分配一组权重，这些权重描述了该点的位置如何由笼顶点决定——笼顶点就像操控木偶的提线，移动任何一个顶点都会带动内部所有点按比例运动。具体而言，GSDeformer 的变形管线分为三个阶段。第一阶段是高斯→代理点云转换：将每个高斯基元转换为显式椭球体，沿三个主轴与椭球面取六个交点作为代理点云 $AP = \{c, p_1^{}+, p_1^{}-, p_2^{}+, p_2^{}-, p_3^{}+\}$，使离散高斯获得可供变形操作的几何锚点。第二阶段是MVC 变形：对代理点云中的每个点，用源笼 $C_s$ 的笼顶点 $\{v_j\}$ 计算其笼坐标 $\{\phi_j\}$（即 $p = \sum_j \phi_j v_j$），然后直接替换为目标笼 $C_d$ 的笼顶点 $\{v'_j\}$，得到变形后位置 $p' = \sum_j \phi_j v'_j$。第三阶段是变换反推与协方差更新：由原始代理点 $AP_s$ 与变形后代理点 $AP_d$ 通过最小二乘拟合 $3 \times 3$ 变换矩阵 $T$，然后将高斯的均值更新为变形后中心，协方差更新为 $\Sigma_d = T \Sigma_s T^{}\top$。这一步至关重要——若不更新协方差，高斯的形状和朝向不会随变形调整，将产生严重的纹理拉伸伪影。对于变形引起的极端弯曲，GSDeformer 引入分裂过程：当变形后高斯沿某轴的代理点弯曲角低于阈值时，将该高斯沿该轴一分为二，每个子高斯的轴长缩放 $k = \sqrt[3]{1/2} \approx 0.79$（近似保持总体积与形状），然后对子高斯重新执行 MVC 变形。这一策略解决了单个高斯无法表示弯曲形状的问题。此外，GSDeformer 还设计了自动化笼生成管线：通过渲染多视角深度图 → T-SDF 重建 → 形态学闭运算 → 约束边折叠抽稀，从已训练的 3DGS 场景中自动构建低面数笼。其中关键创新是两阶段交替抽稀策略——先以约束 QEM 折叠边，再以 MVC 惩罚优化笼顶点位置使负权重趋零——确保笼的几何结构天然适配高质量变形。该方法兼容任何已训练的 vanilla 3DGS 或其变体，无需重训练即可实时编辑。
\textbf{方法信息：} 代码：\url{https://github.com/jhuangbu/gsdeformer-code}。
\textbf{设计动机与边界：} 相对原始 3DGS 以独立高斯和图像重建误差为主的设定，GSDeformer 把问题落在以下可实现的改动上：GSDeformer通过代理点云表示将笼式变形（CBD）应用于3DGS：1）笼构建：自动化生成包围3DGS的三角笼网格。采用「渲染-重建」方法，先从训练好的3DGS渲染多视角图像，用SfM重建网格，再通过约束网格简化（基于QEM的边折叠算法，加入防止表面翻转的约束）生成低多边形笼。2）代理点云生成：从3DGS中的高斯中心位置生成代理点云，为每个代理点计算笼坐标（如均值重心坐标）。3）笼式变形：用户交互式变形笼网格，代理点根据笼坐标计算新位置，然后通过仿射变换将位移传递给对应高斯的位置、旋转和缩放参数。4）分裂处理：当变形引起显著弯曲时，将一个高斯分裂为多个子高斯，每个子高斯对应原高斯的一部分，分别变形以更准确地近似弯曲效果。整个过程不修改3DGS的表示结构，因此可以直接应用于任何已训练的3DGS模型及其变体（如SuGaR、Gaussian Frosting等）。其设计目的不是单独增加一个模块，而是用上述表示、监督或渲染约束消除原流程在目标任务中的具体失配。提出GSDeformer，实现3DGS的笼式变形（cage-based deformation）。通过代理点云表示桥接笼式变形与3DGS，变形应用于点云后转化为高斯变换。引入分裂过程处理变形引起的弯曲。不修改3DGS核心架构，兼容任何已训练的vanilla 3DGS或其变体，无需重训练即可实时编辑，且自动化笼构建。适用时还应保留论文给出的边界：笼质量影响变形效果；自动化笼构建对复杂拓扑可能不理想；极端弯曲时的分裂策略可能产生冗余高斯。
% METHOD: 2405.17811
\paragraph{Mani-GS（三角网格代理操控）（2024.05）}
Mani-GS~\cite{2405_17811} 与 GSDeformer 相邻，但它用三角网格代理而不是 cage。它的目标是支持大变形、局部操控和软体模拟，让用户可以用成熟 mesh 工具或物理模拟来驱动高斯。

流程分为三步。第一步从 3DGS 或神经表面场提取三角 mesh，作为后续编辑的拓扑代理。第二步做 mesh-Gaussian 绑定：每个高斯被关联到一个三角面，其位置、旋转、尺度和颜色与三角形顶点和局部坐标相关，并在训练时用渲染重建损失和形状正则维持外观质量。Shape-aware binding 使高斯尽量扁平贴合三角表面，并允许高斯在三角形间自适应迁移，避免 mesh 提取误差把高斯锁死。第三步才是操控：网格的大变形、局部移动或软体模拟通过绑定关系传递到高斯。相对 GSDeformer 的外部 cage，Mani-GS 的代理更贴近物体表面拓扑，但需要可靠网格作为中介。
\textbf{实现细节：} Mani-GS（2024.05） 提出首个基于三角网格代理的 3DGS 操控框架。直觉上，它将高斯基元"粘"在三角形的局部坐标系上——就像把贴纸贴在三角形瓷砖上，移动或变形瓷砖时贴纸自动跟随，无需逐个重新定位。在此之前，SuGaR 等方法将高斯严格贴附于三角形表面（法线方向缩放趋零），渲染质量高度依赖网格精度——网格不准则高斯无法补偿缺失几何，网格虽精度高但仍有偏差时产生模糊伪影。更宽松的偏移绑定虽然允许高斯离开三角形，但偏移场不随变形迁移，操控后出现渲染扭曲。Mani-GS 的核心改进是三角形感知绑定（Triangle Shape Aware Binding）：不为每个高斯存储全局位置和旋转，而是在每个三角形 $f$ 的局部坐标系中优化局部参数。具体而言，对于由顶点 $(v_1, v_2, v_3)$ 组成的三角形 $f$，Mani-GS 以第一条边方向为第一轴、三角形法向量 $\mathbf{n}^t$ 为第二轴、两者叉积为第三轴，构建局部坐标系 $R^t = [\mathbf{e}_1^t, \mathbf{n}^t, \mathbf{e}_3^t]$。高斯在三角形空间中的局部位置 $\mu^l$ 和局部旋转 $R^l$ 是可优化参数，全局属性通过坐标变换恢复： \begin{equation*} R = R^t R^l, \quad \mu = R^t \mu^l + c^t \end{equation*} 其中 $c^t$ 为三角形中心的全局坐标。这样，当网格变形时三角形局部坐标系 $R^t$ 和中心 $c^t$ 自动更新，高斯的全局位置与旋转无需重新优化即可跟随变化。此外，当三角形在变形中发生伸缩时，高斯的缩放也应按比例调整。Mani-GS 引入自适应向量 $\mathbf{e} = [l_1, \, 0.5(l_1+l_3), \, 0.5(l_2+l_3)]$（其中 $l_i$ 为三角形各边长度），将全局缩放修正为 $s = \mathbf{e} \odot s^l$，确保变形后渲染细节与表面形变一致。该绑定策略对网格精度有较高容忍度（即使使用 Marching Cube 提取的粗糙网格也能保持高保真渲染），并统一支持大变形、局部操控和软体模拟等多种操控类型。
\textbf{方法信息：} Tencent 参与了这篇论文。
\textbf{设计动机与边界：} 相对原始 3DGS 以独立高斯和图像重建误差为主的设定，Mani-GS 把问题落在以下可实现的改动上：Mani-GS通过三角网格代理实现3DGS操控，包含三个阶段。（1）网格提取：从3DGS或神经表面场提取三角网格，支持不同精度和处理时间的提取方法。（2）Mesh-Gaussian绑定：设计三角形感知的绑定方法，将3D高斯绑定到三角形面上。每个高斯的属性（位置、旋转、缩放、颜色）与所在三角形的顶点关联。关键设计是Shape-aware绑定，将高斯分布约束为扁平形状贴合三角形表面，并通过自适应方法允许高斯在三角形间迁移。（3）操控支持：网格变形通过绑定关系直接传递到3D高斯，支持大变形（整体旋转拉伸）、局部操控（移动单个部位）和软体模拟。即使使用从3DGS提取的不精确网格也能获得良好效果。其设计目的不是单独增加一个模块，而是用上述表示、监督或渲染约束消除原流程在目标任务中的具体失配。提出首个基于三角网格代理的3DGS操控框架，通过三角形感知的高斯绑定和自适应方法，实现大变形、局部操控和软体模拟，同时保持操控后高质量渲染。避免为不同类型操控设计不同算法。适用时还应保留论文给出的边界：网格提取质量影响最终效果；复杂拓扑变化（如撕裂）难以处理；绑定过程需要额外优化时间。
% METHOD: 2411.12168
\paragraph{Sketch-guided Cage-based Deformation（草图引导笼式变形）（2024.11）}
Sketch-guided Cage-based 3D Gaussian Splatting Deformation~\cite{2411_12168} 把 cage 变形和 2D 草图先验结合起来。用户只需从一个视角画目标轮廓，系统要把这个 2D 意图转成 3D cage 的平滑变形；直接优化 cage 顶点容易纠缠，且缺少语义合理性。

方法用 cage-based deformation 定义几何映射，高斯的中心和协方差分别通过 cage 变形和局部 Jacobian 更新。为避免 cage 顶点优化不平滑，它引入 Neural Jacobian Fields 来预测或约束局部变形梯度。语义层面，ControlNet 和 SDS 从用户草图中提供 2D 扩散先验，确保形变后的高斯在渲染图上符合草图轮廓和语义。相对 GSDeformer 的手动 cage 拖拽，它把用户输入简化为草图；相对 Drag Your Gaussian，它用 cage/Jacobian 保持几何连续性。
\textbf{实现细节：} 沿笼式路线进一步推进交互性，Sketch-guided Cage（2024.11） 结合笼式变形与 Neural Jacobian Fields 实现细粒度控制，用户通过绘制轮廓草图即可指定变形方向，Neural Jacobian Fields 将草图转化为逐点的 Jacobian 矩阵场，实现比笼顶点拖拽更灵活的局部变形控制。ControlNet+SDS 确保变形结果的语义合理性，支持静态 3DGS 模型的动画化。
\textbf{方法信息：} 代码：\url{https://github.com/tianhaoxie/gs_cage_deform}。
\textbf{设计动机与边界：} 相对原始 3DGS 以独立高斯和图像重建误差为主的设定，Sketch-guided Cage-based 3D Gaussian Splatting Deformation 把问题落在以下可实现的改动上：该方法提出首个草图引导的3DGS变形系统，包含两个核心组件：(1) 几何控制：基于cage-based变形的正则化框架，通过包围网格的顶点移动定义变形函数，将3DGS的质心位置和协方差矩阵分别通过cage变形映射和Jacobian矩阵变换进行变形；引入Neural Jacobian Fields控制cage顶点位置，避免直接优化cage顶点导致的纠缠和不平滑问题；(2) 语义约束：利用ControlNet和SDS(Score Distillation Sampling)从2D扩散先验中获取语义信息，确保变形后的3DGS在语义上合理。用户只需从单一视角绘制轮廓草图即可直观地控制3DGS模型的变形。其设计目的不是单独增加一个模块，而是用上述表示、监督或渲染约束消除原流程在目标任务中的具体失配。首个草图引导的3DGS变形系统，通过cage-based变形与Neural Jacobian Fields结合实现精确细粒度控制，支持静态3DGS模型的动画化适用时还应保留论文给出的边界：论文未明确列出局限，但从单一视角控制可能无法完全定义3D变形；依赖2D扩散先验可能限制变形的多样性。
% METHOD: 2411.18311
\paragraph{Neural Surface Priors（SDF 引导的可编辑高斯）（2024.11）}
Neural Surface Priors for Editable Gaussian Splatting~\cite{2411_18311} 的锚点是用更强几何先验训练可编辑 3DGS。普通高斯不一定贴合真实表面，后续 mesh 编辑传播到高斯时会错位；该方法先训练神经 SDF，再用 SDF 约束高斯分布。

三阶段流程如下。先用 PermutoSDF 从多视图图像得到高质量 SDF，并用 Marching Cubes 提取 mesh。然后进行 SDF-guided 3DGS 训练：距离表面近的高斯获得更高不透明度，固定一个缩放轴为小值使高斯扁平贴合表面，并加入法线对齐正则。训练后，把高斯形状和位置编码到 triangle soup 代理中，每个三角形对应一个高斯或局部高斯集合。用户编辑 mesh 后，根据 mesh 面到 triangle soup 的映射，把变化传播给高斯位置和朝向。相对 Mani-GS，它更强调训练阶段就让高斯服从表面先验。
\textbf{实现细节：} 沿网格代理路线，后续方法从不同角度扩展了能力。Neural Surface Priors（2024.11） 通过 PermutoSDF 获取高质量 SDF 作为几何先验约束高斯分布，设计网格到三角形汤代理的编辑传播机制，支持拓扑感知的编辑操作（如绑定、动画）。
\textbf{实现补充：} 实现上还要核对输入、输出、参数、特征、初始化、训练损失、正则、渲染或推理接口，以及优化中的剪枝、分裂、采样、编码或解码步骤；这些抓手决定该方法到底改了 3DGS 管线的哪一环。
\textbf{方法信息：} 代码：\url{https://github.com/wjakubowska/neuralsurfacepriors}。
\textbf{设计动机与边界：} 相对原始 3DGS 以独立高斯和图像重建误差为主的设定，Neural Surface Priors for Editable Gaussian Splatting 把问题落在以下可实现的改动上：三阶段流程：(1) 神经表面获取：使用PermutoSDF从多视图图像训练神经SDF，独立预处理阶段获得更精细的几何重建。Marching Cubes提取网格供后续编辑。(2) SDF引导的3DGS训练：基于SDF值调节高斯不透明度（距表面近的高斯高不透明度）；固定一个缩放参数为小值使高斯扁平贴合表面；添加正则化对齐高斯法线与SDF预测法线。训练后将高斯形状和位置编码到三角形汤代理——每个三角形对应一个高斯。(3) 编辑传播：将三角形汤代理的每个三角形与最近网格面关联。用户编辑网格后，根据网格面映射关系将修改传播到三角形汤代理，进而更新高斯的位置和朝向。支持绑定、动画等拓扑感知编辑。其设计目的不是单独增加一个模块，而是用上述表示、监督或渲染约束消除原流程在目标任务中的具体失配。提出神经表面先验引导的可编辑3DGS方法。通过PermutoSDF获取高质量SDF作为几何先验约束高斯分布。设计网格到三角形汤代理的编辑传播机制，支持拓扑感知的编辑操作（如绑定、动画），比纯三角形汤方法更灵活。适用时还应保留论文给出的边界：编辑后可能需要重新优化外观；复杂拓扑变化（如切割）不支持；两阶段训练增加总时间
% METHOD: 2510.02034
\paragraph{SemMorph3D（语义感知 3D morphing）（2025.10）}
SemMorph3D~\cite{2510_02034} 处理的是两个或多个 3D 内容之间的语义对应和 morphing。3DGS 的非结构化高斯让直接插值容易碎片化；mesh 有拓扑但可能不够精细。SemMorph3D 的改变是用 mesh-guided Gaussian 策略，把粗 mesh 作为语义和拓扑锚点，高斯承担外观和局部几何细节。

具体方法从多视角图像重建 3DGS，并提取粗基础 mesh。mesh 提供测地线结构和部件连续性，高斯则补偿 mesh 伪影和拓扑限制。双域优化同时包含 geodesic regularization 和 texture-aware constraint：前者作为几何正则项保持形状演化的完整性，后者作为纹理损失保证颜色和纹理在 morphing 过程中连贯；渲染监督则约束中间形态仍能解释输入视图。它不依赖人工标注、类别模板或专门 3D 资产，试图从图像中直接建立无监督语义对应。相对 Mani-GS 的用户操控，它更像自动化的形状与纹理变形。
\textbf{实现细节：} 除了上述两类方法外，还有一些方法从不同角度扩展了 3DGS 变形的能力边界。在输入模态方面，单图像驱动的变形方法（2024.08） 引入位置导数（基于最优传输）捕获长距离运动，设计锚点 ARAP 正则化和粗到细优化策略处理遮挡高斯；视角依赖变形场（2025.04） 允许用户从不同视角对 3DGS 进行 2D 形变，通过平滑插值生成一致的 3D 变形场，支持类似图层编辑的堆叠变形。在网格-高斯混合表示方面，非刚性变形方法（2025.07） 将高斯核直接嵌入显式网格表面，利用网格固有的拓扑和几何先验引导编辑操作；SemMorph3D（2025.10） 通过 mesh 引导的 Gaussian 策略建立无监督语义对应，实现从图像到纹理完整的 3D 形状变形；ObjectMorpher（2026.03） 则将 2D 图像对象提升为可编辑 3DGS，使用 ARAP 约束确保物理合理的形状变化。
\textbf{实现补充：} 实现上还要核对输入、输出、参数、特征、初始化、训练损失、正则、渲染或推理接口，以及优化中的剪枝、分裂、采样、编码或解码步骤；这些抓手决定该方法到底改了 3DGS 管线的哪一环。
\textbf{技术实现补充：} 除了上述两类方法外，还有一些方法从不同角度扩展了 3DGS 变形的能力边界。在输入模态方面，单图像驱动的变形方法（2024.08） 引入位置导数（基于最优传输）捕获长距离运动，设计锚点 ARAP 正则化和粗到细优化策略处理遮挡高斯；视角依赖变形场（2025.04） 允许用户从不同视角对 3DGS 进行 2D 形变，通过平滑插值生成一致的 3D 变形场，支持类似图层编辑的堆叠变形。在网格-高斯混合表示方面，非刚性变形方法（2025.07） 将高斯核直接嵌入显式网格表面，利用网格固有的拓扑和几何先验引导编辑操作；SemMorph3D（2025.10） 通过 mesh 引导的 Gaussian 策略建立无监督语义对应，实现从图像到纹理完整的 3D 形状变形；ObjectMorpher（2026.03） 则将 2D 图像对象提升为可编辑 3DGS，使用 ARAP 约束确保物理合理的形状变化。
\textbf{方法信息：} Tencent 参与了这篇论文。
\textbf{设计动机与边界：} 相对原始 3DGS 以独立高斯和图像重建误差为主的设定，SemMorph3D 把问题落在以下可实现的改动上：SemMorph3D的核心创新在于解决3DGS非结构化特性导致的变形碎片化问题：(1)Mesh-guided Gaussian策略：从重建的3DGS中提取粗略基础mesh作为几何锚点，提供拓扑支架引导非结构化Gaussian的变形，补偿mesh提取的伪影和拓扑限制。(2)双域优化策略：结合测地线正则化(geodesic regularization)保持形状完整性，和纹理感知约束确保颜色演化连贯。两者协同建立无监督语义对应，无需标注数据、专门3D资产或类别特定模板。(3)直接从多视角图像输入生成完全纹理化的拓扑鲁棒3D变形结果，填补了现有方法在2D图像方法和3D几何方法之间的空白。其设计目的不是单独增加一个模块，而是用上述表示、监督或渲染约束消除原流程在目标任务中的具体失配。提出mesh-guided Gaussian策略和双域优化方法，首次实现从图像直接进行无监督语义对应的3D纹理变形，在TexMorph基准上颜色一致性误差(Delta E)降低22.2%，EI降低26.2%。适用时还应保留论文给出的边界：依赖3DGS重建质量；mesh提取可能引入拓扑错误；变形速度较慢

\section{场景修复、物体移除与补全}

修复路线的关键不是“把洞填满”这么简单。3DGS 删除对象后会暴露从未观测过的背景，多视角 2D inpainting 容易在不同视图中生成不一致内容；更进一步，即使视觉上移除了对象，残留高斯或纹理仍可能泄漏语义信息。因此本节方法分别处理未见区域补全、跨视角一致修复、删除残留测量和 2D-3D 联合补全。
% METHOD: 2401.14828
\paragraph{TIP-Editor（编辑、组合与交互修改）（2024.01）}
TIP-Editor~\cite{2401_14828} 关注的问题是：支持 3DGS 场景编辑、组合或交互操控 TIP-Editor 是宽召回补全中发现的论文，主要根据标题、摘要和 arXiv 类别判定其与 3DGS 的关系。
相对原始 3DGS，TIP-Editor 的变化不是只替换损失函数，而是把表示、初始化、约束或网络预测路径重新组织到任务流程中：该工作被召回是因为论文文本包含 Gaussian/3D Gaussian 相关表述；当前分类先确定其任务域和与 3DGS 的相关性，后续证据流水线再补代码、企业和引用信息。
\textbf{实现细节：} 具体实现可以按输入、核心模块和渲染/优化路径来读。输入端首先取得任务所需的图像、相机、文本、时间序列、语义或几何线索，并把这些线索转换成可以约束高斯中心、尺度、旋转、不透明度、颜色或附加特征的条件。核心结构中，方法使用的主要机制是：从摘要层面看，TIP-Editor 围绕“支持 3DGS 场景编辑、组合或交互操控”展开。若其正文或代码后续可读，应进一步核对它具体改动的是表示、初始化、渲染、优化、损失、动态建模、压缩或应用接口中的哪一环。 训练或推理时，这些模块通常先给出初始化或中间特征，再通过可微 Gaussian rasterization、重建损失、感知/语义约束、几何正则或任务损失回传到高斯参数；如果方法包含前馈网络，则网络输出的是高斯属性、变形场、压缩码、语义特征或控制变量，而不是只输出最终 2D 图像。渲染路径仍以显式高斯投影、深度排序和 alpha 合成为基础，只是在采样、属性预测、局部编辑、动态变形、压缩编码或下游任务接口处加入额外约束。这样写入本章时，它与前后方法的关系是：相对原始 3DGS 和本章已有锚点方法，它可以理解为 3DGS extends\_or\_uses 宽召回分类阶段的关系判断，后续需要按论文方法部分复核。
\textbf{证据与边界：} 本轮召回分类尚未逐项核对实验表格；后续证据补全阶段需要读取论文源码/PDF或项目页确认数据集、指标和对比方法。 该条来自宽召回批次；源码下载受 arXiv e-print 网络限制，部分论文暂时只能依据标题、摘要和类别分类，需要在入可信书前做人工/LLM精读复核。 本书只采用本地 paper pack、数据库分类、源码/引用登记和论文摘要中可核验的信息；没有本地证据支撑的速度、排行榜或开源完整性不在正文中强化。\textbf{方法信息：} Tencent 参与；代码：\url{https://github.com/graphdeco-inria/gaussian-splatting}；引用信号：109 citations (3.7/month, domain-relative cutoff 3.0/month, semantic\_scholar)。
\textbf{设计动机与边界：} 相对原始 3DGS 以独立高斯和图像重建误差为主的设定，TIP-Editor 把问题落在以下可实现的改动上：从摘要层面看，TIP-Editor 围绕“支持 3DGS 场景编辑、组合或交互操控”展开。若其正文或代码后续可读，应进一步核对它具体改动的是表示、初始化、渲染、优化、损失、动态建模、压缩或应用接口中的哪一环。其设计目的不是单独增加一个模块，而是用上述表示、监督或渲染约束消除原流程在目标任务中的具体失配。该工作被召回是因为论文文本包含 Gaussian/3D Gaussian 相关表述；当前分类先确定其任务域和与 3DGS 的相关性，后续证据流水线再补代码、企业和引用信息。适用时还应保留论文给出的边界：该条来自宽召回批次；源码下载受 arXiv e-print 网络限制，部分论文暂时只能依据标题、摘要和类别分类，需要在入可信书前做人工/LLM精读复核。
% METHOD: 2401.16663
\paragraph{VR-GS（编辑、组合与交互修改）（2024.01）}
VR-GS~\cite{2401_16663} 关注的问题是：支持 3DGS 场景编辑、组合或交互操控 VR-GS 是宽召回补全中发现的论文，主要根据标题、摘要和 arXiv 类别判定其与 3DGS 的关系。
相对原始 3DGS，VR-GS 的变化不是只替换损失函数，而是把表示、初始化、约束或网络预测路径重新组织到任务流程中：该工作被召回是因为论文文本包含 Gaussian/3D Gaussian 相关表述；当前分类先确定其任务域和与 3DGS 的相关性，后续证据流水线再补代码、企业和引用信息。
\textbf{实现细节：} 具体实现可以按输入、核心模块和渲染/优化路径来读。输入端首先取得任务所需的图像、相机、文本、时间序列、语义或几何线索，并把这些线索转换成可以约束高斯中心、尺度、旋转、不透明度、颜色或附加特征的条件。核心结构中，方法使用的主要机制是：从摘要层面看，VR-GS 围绕“支持 3DGS 场景编辑、组合或交互操控”展开。若其正文或代码后续可读，应进一步核对它具体改动的是表示、初始化、渲染、优化、损失、动态建模、压缩或应用接口中的哪一环。 训练或推理时，这些模块通常先给出初始化或中间特征，再通过可微 Gaussian rasterization、重建损失、感知/语义约束、几何正则或任务损失回传到高斯参数；如果方法包含前馈网络，则网络输出的是高斯属性、变形场、压缩码、语义特征或控制变量，而不是只输出最终 2D 图像。渲染路径仍以显式高斯投影、深度排序和 alpha 合成为基础，只是在采样、属性预测、局部编辑、动态变形、压缩编码或下游任务接口处加入额外约束。这样写入本章时，它与前后方法的关系是：相对原始 3DGS 和本章已有锚点方法，它可以理解为 3DGS extends\_or\_uses 宽召回分类阶段的关系判断，后续需要按论文方法部分复核。
\textbf{证据与边界：} 本轮召回分类尚未逐项核对实验表格；后续证据补全阶段需要读取论文源码/PDF或项目页确认数据集、指标和对比方法。 该条来自宽召回批次；源码下载受 arXiv e-print 网络限制，部分论文暂时只能依据标题、摘要和类别分类，需要在入可信书前做人工/LLM精读复核。 本书只采用本地 paper pack、数据库分类、源码/引用登记和论文摘要中可核验的信息；没有本地证据支撑的速度、排行榜或开源完整性不在正文中强化。\textbf{方法信息：} Amazon、Project page lists Style3D Research and Amazon affiliations 参与；引用信号：265 citations (9.0/month, domain-relative cutoff 3.0/month, semantic\_scholar)。
\textbf{设计动机与边界：} 相对原始 3DGS 以独立高斯和图像重建误差为主的设定，VR-GS 把问题落在以下可实现的改动上：从摘要层面看，VR-GS 围绕“支持 3DGS 场景编辑、组合或交互操控”展开。若其正文或代码后续可读，应进一步核对它具体改动的是表示、初始化、渲染、优化、损失、动态建模、压缩或应用接口中的哪一环。其设计目的不是单独增加一个模块，而是用上述表示、监督或渲染约束消除原流程在目标任务中的具体失配。该工作被召回是因为论文文本包含 Gaussian/3D Gaussian 相关表述；当前分类先确定其任务域和与 3DGS 的相关性，后续证据流水线再补代码、企业和引用信息。适用时还应保留论文给出的边界：该条来自宽召回批次；源码下载受 arXiv e-print 网络限制，部分论文暂时只能依据标题、摘要和类别分类，需要在入可信书前做人工/LLM精读复核。
% METHOD: 2402.01459
\paragraph{GaMeS（编辑、组合与交互修改）（2024.02）}
GaMeS~\cite{2402_01459} 关注的问题是：支持 3DGS 场景编辑、组合或交互操控 GaMeS 是宽召回补全中发现的论文，主要根据标题、摘要和 arXiv 类别判定其与 3DGS 的关系。
相对原始 3DGS，GaMeS 的变化不是只替换损失函数，而是把表示、初始化、约束或网络预测路径重新组织到任务流程中：该工作被召回是因为论文文本包含 Gaussian/3D Gaussian 相关表述；当前分类先确定其任务域和与 3DGS 的相关性，后续证据流水线再补代码、企业和引用信息。
\textbf{实现细节：} 具体实现可以按输入、核心模块和渲染/优化路径来读。输入端首先取得任务所需的图像、相机、文本、时间序列、语义或几何线索，并把这些线索转换成可以约束高斯中心、尺度、旋转、不透明度、颜色或附加特征的条件。核心结构中，方法使用的主要机制是：从摘要层面看，GaMeS 围绕“支持 3DGS 场景编辑、组合或交互操控”展开。若其正文或代码后续可读，应进一步核对它具体改动的是表示、初始化、渲染、优化、损失、动态建模、压缩或应用接口中的哪一环。 训练或推理时，这些模块通常先给出初始化或中间特征，再通过可微 Gaussian rasterization、重建损失、感知/语义约束、几何正则或任务损失回传到高斯参数；如果方法包含前馈网络，则网络输出的是高斯属性、变形场、压缩码、语义特征或控制变量，而不是只输出最终 2D 图像。渲染路径仍以显式高斯投影、深度排序和 alpha 合成为基础，只是在采样、属性预测、局部编辑、动态变形、压缩编码或下游任务接口处加入额外约束。这样写入本章时，它与前后方法的关系是：相对原始 3DGS 和本章已有锚点方法，它可以理解为 3DGS extends\_or\_uses 宽召回分类阶段的关系判断，后续需要按论文方法部分复核。
\textbf{证据与边界：} 本轮召回分类尚未逐项核对实验表格；后续证据补全阶段需要读取论文源码/PDF或项目页确认数据集、指标和对比方法。 该条来自宽召回批次；源码下载受 arXiv e-print 网络限制，部分论文暂时只能依据标题、摘要和类别分类，需要在入可信书前做人工/LLM精读复核。 本书只采用本地 paper pack、数据库分类、源码/引用登记和论文摘要中可核验的信息；没有本地证据支撑的速度、排行榜或开源完整性不在正文中强化。\textbf{方法信息：} 代码：\url{https://github.com/waczjoan/gaussian-mesh-splatting}；引用信号：88 citations (3.0/month, domain-relative cutoff 3.0/month, semantic\_scholar)；代码关注：waczjoan/gaussian-mesh-splatting has 505 stars (17.3/month, domain-relative cutoff 9.8/month)。
\textbf{设计动机与边界：} 相对原始 3DGS 以独立高斯和图像重建误差为主的设定，GaMeS 把问题落在以下可实现的改动上：从摘要层面看，GaMeS 围绕“支持 3DGS 场景编辑、组合或交互操控”展开。若其正文或代码后续可读，应进一步核对它具体改动的是表示、初始化、渲染、优化、损失、动态建模、压缩或应用接口中的哪一环。其设计目的不是单独增加一个模块，而是用上述表示、监督或渲染约束消除原流程在目标任务中的具体失配。该工作被召回是因为论文文本包含 Gaussian/3D Gaussian 相关表述；当前分类先确定其任务域和与 3DGS 的相关性，后续证据流水线再补代码、企业和引用信息。适用时还应保留论文给出的边界：该条来自宽召回批次；源码下载受 arXiv e-print 网络限制，部分论文暂时只能依据标题、摘要和类别分类，需要在入可信书前做人工/LLM精读复核。
% METHOD: 2403.11868
\paragraph{View-Consistent 3D Editing with Gaussian Splatting（编辑、组合与交互修改）（2024.03）}
View-Consistent 3D Editing with Gaussian Splatting~\cite{2403_11868} 关注的问题是：支持 3DGS 场景编辑、组合或交互操控 View-Consistent 3D Editing with Gaussian Splatting 是宽召回补全中发现的论文，主要根据标题、摘要和 arXiv 类别判定其与 3DGS 的关系。
相对原始 3DGS，View-Consistent 3D Editing with Gaussian Splatting 的变化不是只替换损失函数，而是把表示、初始化、约束或网络预测路径重新组织到任务流程中：该工作被召回是因为论文文本包含 Gaussian/3D Gaussian 相关表述；当前分类先确定其任务域和与 3DGS 的相关性，后续证据流水线再补代码、企业和引用信息。
\textbf{实现细节：} 具体实现可以按输入、核心模块和渲染/优化路径来读。输入端首先取得任务所需的图像、相机、文本、时间序列、语义或几何线索，并把这些线索转换成可以约束高斯中心、尺度、旋转、不透明度、颜色或附加特征的条件。核心结构中，方法使用的主要机制是：从摘要层面看，View-Consistent 3D Editing with Gaussian Splatting 围绕“支持 3DGS 场景编辑、组合或交互操控”展开。若其正文或代码后续可读，应进一步核对它具体改动的是表示、初始化、渲染、优化、损失、动态建模、压缩或应用接口中的哪一环。 训练或推理时，这些模块通常先给出初始化或中间特征，再通过可微 Gaussian rasterization、重建损失、感知/语义约束、几何正则或任务损失回传到高斯参数；如果方法包含前馈网络，则网络输出的是高斯属性、变形场、压缩码、语义特征或控制变量，而不是只输出最终 2D 图像。渲染路径仍以显式高斯投影、深度排序和 alpha 合成为基础，只是在采样、属性预测、局部编辑、动态变形、压缩编码或下游任务接口处加入额外约束。这样写入本章时，它与前后方法的关系是：相对原始 3DGS 和本章已有锚点方法，它可以理解为 3DGS extends\_or\_uses 宽召回分类阶段的关系判断，后续需要按论文方法部分复核。
\textbf{证据与边界：} 本轮召回分类尚未逐项核对实验表格；后续证据补全阶段需要读取论文源码/PDF或项目页确认数据集、指标和对比方法。 该条来自宽召回批次；源码下载受 arXiv e-print 网络限制，部分论文暂时只能依据标题、摘要和类别分类，需要在入可信书前做人工/LLM精读复核。 本书只采用本地 paper pack、数据库分类、源码/引用登记和论文摘要中可核验的信息；没有本地证据支撑的速度、排行榜或开源完整性不在正文中强化。\textbf{方法信息：} 代码：\url{https://github.com/yuxuan-w/vcedit}；引用信号：85 citations (3.1/month, domain-relative cutoff 3.0/month, semantic\_scholar)。
\textbf{设计动机与边界：} 相对原始 3DGS 以独立高斯和图像重建误差为主的设定，View-Consistent 3D Editing with Gaussian Splatting 把问题落在以下可实现的改动上：从摘要层面看，View-Consistent 3D Editing with Gaussian Splatting 围绕“支持 3DGS 场景编辑、组合或交互操控”展开。若其正文或代码后续可读，应进一步核对它具体改动的是表示、初始化、渲染、优化、损失、动态建模、压缩或应用接口中的哪一环。其设计目的不是单独增加一个模块，而是用上述表示、监督或渲染约束消除原流程在目标任务中的具体失配。该工作被召回是因为论文文本包含 Gaussian/3D Gaussian 相关表述；当前分类先确定其任务域和与 3DGS 的相关性，后续证据流水线再补代码、企业和引用信息。适用时还应保留论文给出的边界：该条来自宽召回批次；源码下载受 arXiv e-print 网络限制，部分论文暂时只能依据标题、摘要和类别分类，需要在入可信书前做人工/LLM精读复核。
% METHOD: 2411.19233
\paragraph{Gaussians-to-Life（文本驱动静态场景动画化）（2024）}
Gaussians-to-Life~\cite{2411_19233} 把编辑对象从静态外观和几何扩展到运动。输入是已有高质量静态 3DGS 场景和文本提示，例如让某个对象动起来；输出是带 3D 运动的高斯场景。它不依赖角色骨骼或特定类别模板，因此适合复杂重建场景中的多种物体。

方法先从某个视角渲染场景，并用视频扩散模型根据文本提示生成动画视频。然后通过鲁棒的 2D-to-3D motion lifting，把视频中的 2D 运动估计成 3D 空间中高斯基元的位移或变换。实现上它接入 threestudio 和 DynamiCrafter：先准备静态 3DGS、文本提示和初始视频，再用配置文件驱动训练，把视频中的运动提升为选定高斯集合的时间变化。最后系统把运动应用到高斯场，使场景在自由视角下呈现一致动画。相对 4DStyleGaussian 的动态场景风格化，Gaussians-to-Life 的目标是把静态 3DGS 激活动起来；相对 Drag Your Gaussian，它的运动来源是文本驱动视频先验而非用户拖拽点。
\textbf{设计动机与边界：} 相对原始 3DGS 以独立高斯和图像重建误差为主的设定，Gaussians-to-Life 把问题落在以下可实现的改动上：Gaussians2Life方法包含三个关键步骤：首先使用强大的视频扩散模型根据文本提示从3D场景的特定视角生成动画视频；然后通过鲁棒的2D到3D运动提升技术，将视频中的2D运动转换为3D空间中高斯基元的位移；最后对选定的3DGS基元施加运动变换实现动画。该方法不依赖先验的角色模型，可处理各种物体类别，并通过3DGS的显式表示实现高效渲染。其设计目的不是单独增加一个模块，而是用上述表示、监督或渲染约束消除原流程在目标任务中的具体失配。提出首个能对任意预存在的高质量3DGS场景进行文本驱动动画的方法，不限于角色动画或单个物体，支持复杂场景中多种物体类别的动画化，创建一致且沉浸式的3D体验。适用时还应保留论文给出的边界：动画质量受限于视频扩散模型生成的运动质量；复杂运动可能导致多视角不一致；对细粒度运动控制能力有限。
\textbf{方法信息：} Google Research 参与了这篇论文；代码：\url{https://github.com/wimmerth/gaussians2life}。

% METHOD: 2502.05176
\paragraph{AuraFusion360（360 度无界场景参考式修复）（2025.02）}
AuraFusion360~\cite{2502_05176} 面向 360 度无界场景的对象移除与背景补全。普通 3D inpainting 多在前向场景或有限视角中评估；360 度场景中，相机可能从任意方向看向被删除区域，未见背景的几何和外观都更难对齐。

方法先训练带对象掩码的高斯表示并移除目标对象。随后做 depth-aware unseen mask generation：利用深度变形估计各视角未见区域轮廓，把边界框提示给 SAM2，生成精确未见掩码，区分可见背景和真正需要修复的区域。Adaptive Guided Depth Diffusion 从参考视角反投影对齐点，初始化未见区域高斯，并引导 RGB 修复，无需额外训练。最后用 SDEdit 对多视角修复图像做细节增强。相对 3DGIC，它更专注 360 度无界场景和大视角变化。
\textbf{实现细节：} 沿多视角扩散修复路线，后续方法将能力扩展到更复杂的场景类型。AuraFusion360（2025.02） 是首个针对 360 度无界场景的参考式 3D 修复方法，通过深度感知未见掩码生成、自适应引导深度扩散和 SDEdit 细节增强实现高质量对象移除。LightHarmony3D（2026.03） 专门解决物体插入中的光照一致性问题，通过生成式 HDR 环境图预测和 PBR 引导的阴影合成确保物理一致的阴影整合。
\textbf{实现补充：} 为避免只把它读成名称或数据点，还应把实现接口明确拆开：输入是什么，输出是高斯参数、图像、轨迹、标注还是评测指标，哪些参数、网络模块、训练损失或渲染接口真正接到原始 3DGS / Scaffold-GS / 4D-GS 管线上，以及剪枝、分裂、采样、解码、优化或评测协议在哪一步发生。若该条目是survey、benchmark 或 dataset，这里的“实现”对应论文集合、任务划分、相机/轨迹/标注组织、评价指标和基线配置，而不是虚构新的 rasterizer；若是算法条目，则对应高斯位置、尺度、旋转、不透明度、特征或 MLP 解码如何进入训练和渲染流程。
\textbf{方法信息：} NVIDIA 参与了这篇论文；代码：\url{https://github.com/kkennethwu/AuraFusion360_official}。



% METHOD: 2502.11801
\paragraph{3DGIC（深度引导跨视角一致修复）（2025.02）}
3D Gaussian Inpainting with Depth-Guided Cross-View Consistency~\cite{2502_11801} 的核心是用深度把不同视角的可见背景联系起来。纯 2D 修复模型在每张图上都能补洞，但补出的背景可能不是同一个 3D 表面；3DGIC 的改变是在修复前先推断哪些背景像素可从其他视角投影过来。

第一阶段用 3DGS 渲染深度图，把其他视角中可见的背景像素投影到目标视图，据此更新修复掩码，只让真正未知区域交给 2D inpainting。第二阶段把修复模型输出作为监督更新 3DGS，并通过跨视角背景对应保持一致。这样，已有真实背景优先被复用，生成模型只补不可见部分。相对 AuraFusion360，它的技术焦点是 depth-guided mask refinement 和 cross-view consistency，而不是参考式 360 度数据集。
\textbf{实现细节：} GScream 的交叉注意力特征传播虽然在特征空间建立了修复区域与可见区域的关联，但依赖 Scaffold-GS 的锚点结构，且传播效果受限于锚点采样的覆盖质量。3DGIC（2025.02） 从这一局限出发，其关键观察是：GScream 的特征传播在修复区域边界处容易产生模糊，因为注意力机制难以精确区分"真正需要修复的像素"和"从其他视角已可见的背景像素"。为此，3DGIC 放弃了特征空间的传播策略，转而在几何空间利用渲染深度做跨视角可见性分析。具体地，对参考视角的掩码区域 $M_1$，3DGIC 将其余视角的可见背景像素通过渲染深度 $D_{2:K}$ 投影到 3D 空间，再反投影回参考视角，提取落在 $M_1$ 内的像素坐标——这些像素在其他视角中已可见，无需修复。通过确定性几何投影而非特征匹配，3DGIC 将原始掩码 $M_1$ 精确缩减为 $M'_1$，仅保留真正被所有视角遮挡的区域。
\textbf{实现补充：} 实现上还要核对输入、输出、参数、特征、初始化、训练损失、正则、渲染或推理接口，以及优化中的剪枝、分裂、采样、编码或解码步骤；这些抓手决定该方法到底改了 3DGS 管线的哪一环。
\textbf{方法信息：} NVIDIA 参与了这篇论文。
\textbf{设计动机与边界：} 相对原始 3DGS 以独立高斯和图像重建误差为主的设定，3D Gaussian Inpainting with Depth-Guided Cross-View Consistency 把问题落在以下可实现的改动上：3DGIC包含两个阶段：(1)深度引导修复掩码推断：给定多视角图像和物体掩码，利用3DGS渲染的深度图将其他视角的可见背景像素投影到目标视角，推断精确的修复掩码，确保仅使用跨视角一致的背景信息；(2)修复引导的3DGS精化：基于精化掩码更新3DGS，使用2D修复模型的输出作为监督进行多视角一致的3D修复。通过将不同视角的背景像素通过3D投影关联，解决了纯2D修复方法的多视角不一致问题。其设计目的不是单独增加一个模块，而是用上述表示、监督或渲染约束消除原流程在目标任务中的具体失配。提出3DGIC框架，利用渲染深度信息引导跨视角一致的3D高斯修复。通过深度引导的修复掩码更新策略，利用跨视角可见的背景像素精化3DGS修复结果，在基准数据集上定量和定性超越当前SOTA 3D修复方法。适用时还应保留论文给出的边界：依赖2D修复模型的质量；深度渲染误差可能影响掩码精度；对大型被遮挡区域的修复效果有限。
% METHOD: 2503.17574
\paragraph{Semantic Residuals after Object Removal（删除后的语义残留测量）（2025.03）}
Is there anything left?~\cite{2503_17574} 不提出新的主流编辑器，而是提出一个很重要的审计问题：物体从 3DGS 中“看起来”删除了，是否还可被语义模型推断出来？这对隐私、安全和数据发布很关键，因为残余高斯或纹理可能仍保留目标类别线索。

评估流程在物体移除后的区域运行开放词汇分割或识别模型，搜索是否仍能识别被删除对象，并用置信度、IoU 下降和深度变化等指标量化 semantic residual。作者还提出空间一致性和语义一致性 refinement：在初始删除后，继续查找可能残留的高斯基元，结合 SAM 掩码、CLIP 语义特征和多视角几何约束进一步清理；这一步相当于把“删除是否干净”变成可迭代的检查和优化问题。相对 3DGIC/AuraFusion360 关注补全质量，它关注删除是否真的干净，是编辑章节中少见的质量和隐私审计方法。
\textbf{实现细节：} 在质量评估方面，Remove360（2025.08） 和语义残差测量（2025.03） 首次系统研究 3DGS 中物体移除后的语义残留问题，揭示几何移除与语义擦除之间的显著差距。高保真修复方法（2025.07） 通过自动掩码精炼和区域级不确定性引导优化实现多视角一致且照片级真实的修复。DeclutterNeRF（2025.04） 通过联合多视角可学习相机参数优化和遮挡退火正则化，无需生成先验即可移除复杂遮挡。
\textbf{方法信息：} 代码已公开，仓库为 \texttt{simonakocour/anything\_left} (4 stars)。入选依据：企业:否；代码:是；星标:否；引用:否。
\textbf{设计动机与边界：} 相对原始 3DGS 以独立高斯和图像重建误差为主的设定，Is there anything left? Measuring semantic residuals of objects removed from 3D Gaussian S 把问题落在以下可实现的改动上：本文提出了一种定量评估框架，用于测量3D场景（基于3DGS和NeRF表示）中物体被移除后是否仍有可被推理的语义残差。核心思想是：如果移除后仍能通过语义模型识别出被移除的物体类别，则场景不安全。评估流程包括：(1)使用开放词汇分割模型在被移除区域搜索残余语义信息；(2)通过置信度分数量化残差程度。此外，提出基于空间一致性和语义一致性的精炼移除方法，在初始移除操作后进一步清理可能遗漏的高斯基元。方法结合了SAM分割和CLIP语义特征，确保移除的完整性。其设计目的不是单独增加一个模块，而是用上述表示、监督或渲染约束消除原流程在目标任务中的具体失配。首次提出定量评估3D场景中物体移除后语义残差的方法，通过度量移除操作后是否仍可推理出被移除内容来判断隐私安全性，并提出基于空间和语义一致性的移除精炼方法。适用时还应保留论文给出的边界：目前仅评估了有限数量的移除方法和场景类型；语义残差的定义可能不涵盖所有隐私泄露方式。
% METHOD: 2511.20924
\paragraph{GaINeR（编辑、组合与交互修改）（2025.11）}
GaINeR~\cite{2511_20924} 关注的问题是：支持 3DGS 场景编辑、组合或交互操控 GaINeR 是宽召回补全中发现的论文，主要根据标题、摘要和 arXiv 类别判定其与 3DGS 的关系。
作为 supplementary 类条目，GaINeR 不应被写成单纯的新渲染算子。它更像是把 3DGS 放进 场景编辑与操控 的协议、系统或应用边界里：该工作被召回是因为论文文本包含 Gaussian/3D Gaussian 相关表述；当前分类先确定其任务域和与 3DGS 的相关性，后续证据流水线再补代码、企业和引用信息。
\textbf{实现细节：} 具体实现可以按输入、核心模块和渲染/优化路径来读。输入端首先取得任务所需的图像、相机、文本、时间序列、语义或几何线索，并把这些线索转换成可以约束高斯中心、尺度、旋转、不透明度、颜色或附加特征的条件。核心结构中，方法使用的主要机制是：从摘要层面看，GaINeR 围绕“支持 3DGS 场景编辑、组合或交互操控”展开。若其正文或代码后续可读，应进一步核对它具体改动的是表示、初始化、渲染、优化、损失、动态建模、压缩或应用接口中的哪一环。 训练或推理时，这些模块通常先给出初始化或中间特征，再通过可微 Gaussian rasterization、重建损失、感知/语义约束、几何正则或任务损失回传到高斯参数；如果方法包含前馈网络，则网络输出的是高斯属性、变形场、压缩码、语义特征或控制变量，而不是只输出最终 2D 图像。渲染路径仍以显式高斯投影、深度排序和 alpha 合成为基础，只是在采样、属性预测、局部编辑、动态变形、压缩编码或下游任务接口处加入额外约束。这样写入本章时，它与前后方法的关系是：相对原始 3DGS 和本章已有锚点方法，它可以理解为 3DGS extends\_or\_uses 宽召回分类阶段的关系判断，后续需要按论文方法部分复核。
\textbf{证据与边界：} 本轮召回分类尚未逐项核对实验表格；后续证据补全阶段需要读取论文源码/PDF或项目页确认数据集、指标和对比方法。 该条来自宽召回批次；源码下载受 arXiv e-print 网络限制，部分论文暂时只能依据标题、摘要和类别分类，需要在入可信书前做人工/LLM精读复核。 本书只采用本地 paper pack、数据库分类、源码/引用登记和论文摘要中可核验的信息；没有本地证据支撑的速度、排行榜或开源完整性不在正文中强化。\textbf{方法信息：} 代码：\url{https://github.com/wjakubowska/gainer}。
\textbf{设计动机与边界：} 相对原始 3DGS 以独立高斯和图像重建误差为主的设定，GaINeR 把问题落在以下可实现的改动上：从摘要层面看，GaINeR 围绕“支持 3DGS 场景编辑、组合或交互操控”展开。若其正文或代码后续可读，应进一步核对它具体改动的是表示、初始化、渲染、优化、损失、动态建模、压缩或应用接口中的哪一环。其设计目的不是单独增加一个模块，而是用上述表示、监督或渲染约束消除原流程在目标任务中的具体失配。该工作被召回是因为论文文本包含 Gaussian/3D Gaussian 相关表述；当前分类先确定其任务域和与 3DGS 的相关性，后续证据流水线再补代码、企业和引用信息。适用时还应保留论文给出的边界：该条来自宽召回批次；源码下载受 arXiv e-print 网络限制，部分论文暂时只能依据标题、摘要和类别分类，需要在入可信书前做人工/LLM精读复核。
% METHOD: 2606.27584
\paragraph{CoIn（2D-3D 综合修复与高斯引导）（2026）}
CoIn~\cite{2606_27584} 处理的是更一般的 3D 场景修复：不只删除物体，也包括任意形状区域补全和对象插入。许多 3DGS inpainting 方法要求用户给出精确多视图 mask，或者只在物体移除后补背景；CoIn 的改变是让 2D 扩散修复和 3D 高斯重建来回协作，用 2D 模型负责生成内容，用 3DGS 负责几何对齐和跨视角一致。

流程可以理解为 2D-to-3D、3D-to-2D、再 2D-to-3D 的闭环。第一阶段先用 2D 扩散模型在参考视图上根据任意形状 mask 生成初始修复结果，再通过 Reference Adaptive GS 重建粗 3D 场景；其中 Feature Attention 用来强调修复区域相关的纹理和语义，使粗重建不要被未修复区域主导。第二阶段利用高斯场景提供几何引导：从当前 3DGS 渲染出的深度、特征或参考信息被 warping 到目标视图，作为扩散过程的 reference feature guidance，并通过 consistency loss guidance 抑制不同视角各自生成、互相不一致的问题。第三阶段再把修复后的多视图结果优化回 3DGS，并加入 texture-enhancing discriminator，补强高频纹理和真实感。相对 3DGIC 优先复用已有背景像素的深度投影路线，CoIn 更依赖生成模型和高斯几何之间的双向反馈，覆盖范围也更宽。
\textbf{方法信息：} 代码：\url{https://github.com/hana001106/CoIn}；本地检查时仓库 README 仍显示实现整理中，后续需要继续确认完整代码释放。
\textbf{设计动机与边界：} 相对原始 3DGS 以独立高斯和图像重建误差为主的设定，CoIn 把问题落在以下可实现的改动上：先由 2D 扩散模型生成初始修复图像，再用 Reference Adaptive GS 和 Feature Attention 重建粗 3D 场景，并通过多阶段一致性优化减少跨视图冲突。其设计目的不是单独增加一个模块，而是用上述表示、监督或渲染约束消除原流程在目标任务中的具体失配。突破现有 GS 编辑方法依赖精确多视图 mask 且多用于物体移除的限制，支持对象插入等多样任务。适用时还应保留论文给出的边界：2D 扩散先验可能造成几何不一致或语义幻觉；多阶段优化成本和复杂遮挡区域仍是挑战。
% METHOD: 2607.05347
\paragraph{WildSplat（编辑、组合与交互修改）（2026.07）}
WildSplat~\cite{2607_05347} 关注的问题是：支持 3DGS 场景编辑、组合或交互操控 WildSplat 是宽召回补全中发现的论文，主要根据标题、摘要和 arXiv 类别判定其与 3DGS 的关系。
相对原始 3DGS，WildSplat 的变化不是只替换损失函数，而是把表示、初始化、约束或网络预测路径重新组织到任务流程中：该工作被召回是因为论文文本包含 Gaussian/3D Gaussian 相关表述；当前分类先确定其任务域和与 3DGS 的相关性，后续证据流水线再补代码、企业和引用信息。
\textbf{实现细节：} 具体实现可以按输入、核心模块和渲染/优化路径来读。输入端首先取得任务所需的图像、相机、文本、时间序列、语义或几何线索，并把这些线索转换成可以约束高斯中心、尺度、旋转、不透明度、颜色或附加特征的条件。核心结构中，方法使用的主要机制是：从摘要层面看，WildSplat 围绕“支持 3DGS 场景编辑、组合或交互操控”展开。若其正文或代码后续可读，应进一步核对它具体改动的是表示、初始化、渲染、优化、损失、动态建模、压缩或应用接口中的哪一环。 训练或推理时，这些模块通常先给出初始化或中间特征，再通过可微 Gaussian rasterization、重建损失、感知/语义约束、几何正则或任务损失回传到高斯参数；如果方法包含前馈网络，则网络输出的是高斯属性、变形场、压缩码、语义特征或控制变量，而不是只输出最终 2D 图像。渲染路径仍以显式高斯投影、深度排序和 alpha 合成为基础，只是在采样、属性预测、局部编辑、动态变形、压缩编码或下游任务接口处加入额外约束。这样写入本章时，它与前后方法的关系是：相对原始 3DGS 和本章已有锚点方法，它可以理解为 3DGS extends\_or\_uses 宽召回分类阶段的关系判断，后续需要按论文方法部分复核。
\textbf{证据与边界：} 本轮召回分类尚未逐项核对实验表格；后续证据补全阶段需要读取论文源码/PDF或项目页确认数据集、指标和对比方法。 该条来自宽召回批次；源码下载受 arXiv e-print 网络限制，部分论文暂时只能依据标题、摘要和类别分类，需要在入可信书前做人工/LLM精读复核。 本书只采用本地 paper pack、数据库分类、源码/引用登记和论文摘要中可核验的信息；没有本地证据支撑的速度、排行榜或开源完整性不在正文中强化。\textbf{方法信息：} 代码：\url{https://github.com/yfujimura/wildsplatter}。
\textbf{设计动机与边界：} 相对原始 3DGS 以独立高斯和图像重建误差为主的设定，WildSplat 把问题落在以下可实现的改动上：从摘要层面看，WildSplat 围绕“支持 3DGS 场景编辑、组合或交互操控”展开。若其正文或代码后续可读，应进一步核对它具体改动的是表示、初始化、渲染、优化、损失、动态建模、压缩或应用接口中的哪一环。其设计目的不是单独增加一个模块，而是用上述表示、监督或渲染约束消除原流程在目标任务中的具体失配。该工作被召回是因为论文文本包含 Gaussian/3D Gaussian 相关表述；当前分类先确定其任务域和与 3DGS 的相关性，后续证据流水线再补代码、企业和引用信息。适用时还应保留论文给出的边界：该条来自宽召回批次；源码下载受 arXiv e-print 网络限制，部分论文暂时只能依据标题、摘要和类别分类，需要在入可信书前做人工/LLM精读复核。
